% ======================================================================
% Chapter 17: Quantum Chromodynamics
% (book pages 545-597; PDF pages 568-621)
% ======================================================================
\chapter{量子色动力学}

构造强相互作用的现实模型的关键是渐近自由现象。第~14 章描述了这一现象的实验发现，而第~16 章给出了非阿贝尔规范理论是渐近自由的这一理论证明。我们现在已经准备好探索这些发现的后果。

我们将首先论证，强相互作用模型最自然的候选者是规范群为 $SU(3)$、与基础表示中的费米子（夸克）耦合的非阿贝尔规范理论。这个理论被称为量子色动力学，即 QCD。在第~17.1 节对 QCD 作了一些一般性讨论之后，我们将在第~17.2 节到第~17.4 节研究若干具体的 QCD 散射过程。然而，QCD 最令人感兴趣的应用具有某种更为精细的性质；它在于预言第~14 章中讨论的 Bjorken 标度关系缓慢破坏的模式。第~17.5 节将发展理解这些破坏所需的额外理论工具。

尽管本章包含许多对实验的引用，但我们提醒读者，对于 QCD 正如对于 QED 或临界现象一样，本书主要是理论方法的教科书，而不是实验数据的综述与解读。强相互作用物理实验技术与结果的细节在若干优秀文献中都有综述（参见参考文献目录）。我们希望本章能给出阐明和解读这些结果所必需的理论基础。

\section{从夸克到 QCD}

我们当前对强相互作用的理论图景始于对构成质子及其他强子的基本费米子的辨认。随着人们对这些费米子的性质理解得越来越清楚，它们相互作用的性质也受到越来越强的约束，其方式最终导致一个唯一的候选理论。为了理解这个理论的唯一性，我们以这个理论如何产生的简化历史来开始本章。

1963 年，Gell-Mann 和 Zweig 提出了一个模型，用称为\emph{夸克}的基本组分解释了强相互作用粒子的谱。介子被认为是夸克-反夸克束缚态。确实，最轻的介子正好具有证实这种解释所需的量子数；它们是奇宇称的自旋-0 和自旋-1 态，正如我们在第~3 章中对零轨道角动量的费米子-反费米子束缚态所发现的那样。重子被解释为三个夸克的束缚态。为了解释强子的电荷和其他量子数，Gell-Mann 和 Zweig 需要假设三种夸克：上（$u$）、下（$d$）和奇异（$s$）。此后发现的其他强子要求再存在三种夸克：粲（$c$）、底（$b$）和顶（$t$）。为了使重子具有整数电荷，需要给夸克分配分数电荷：$u$、$c$、$t$ 为 $+2/3$，$d$、$s$、$b$ 为 $-1/3$。于是，例如，质子是 $uud$ 的束缚态，而中子是 $udd$ 的束缚态。这六种夸克按惯例称为\emph{味}。

夸克模型在预言新的强子态、以及解释不同强子之间电磁和弱相互作用跃迁的强度方面取得了巨大成功。特别是，夸克模型自然地包含了强相互作用粒子之间最重要的对称关系。如果假设 $u$ 和 $d$ 夸克具有相同的质量和相互作用，那么作为 $u$ 和 $d$ 态幺正旋转的 $SU(2)$ 群，
\begin{equation}
  \begin{pmatrix}u\\ d\end{pmatrix} \to U\begin{pmatrix}u\\ d\end{pmatrix},
  \label{eq:17.1}
\end{equation}
应该是强相互作用的对称性。确实，无论在原子核物理还是基本粒子物理中，量子态都构成这种 $SU(2)$ 对称性的多重态，称为\emph{同位旋}。类似地，由于奇异夸克只比 $u$ 和 $d$ 夸克稍重，考虑三重态 $(u, d, s)$ 幺正变换的对称性是有意义的。Gell-Mann 和 Ne'eman 证明了基本粒子自然地填满这种 $SU(3)$ 对称性的不可约表示。

尽管原始夸克模型在唯象上取得了成功，它仍然有两个严重问题。第一，尽管付出了相当大的努力，仍找不到带分数电荷的自由粒子。第二，重子谱要求假设三个夸克的波函数在夸克自旋和味量子数的交换下完全对称，这与自旋必须为 $1/2$ 的夸克应服从 Fermi-Dirac 统计的预期相矛盾。对这种对称性的需要最清楚地体现在如下事实中：核子的一个最轻激发态是电荷为 $+2$ 的自旋-$3/2$ 粒子 $\Delta^{++}$。这个粒子很容易被解释为轨道角动量为零且三个夸克自旋平行的 $uuu$ 束缚态。

为了使重子谱与自旋-统计定理相容，Han 和 Nambu、Greenberg 以及 Gell-Mann 提出夸克携带一个额外的、未被观测到的量子数，称为\emph{色}。他们引入了这样一个特设假设：重子波函数在色量子数上必须完全反对称。于是，如果夸克波函数在自旋和味上完全对称，那么它们在总体上完全反对称，与 Fermi-Dirac 统计一致。最简单的色模型是把夸克指定到一个新的内部 $SU(3)$ 整体对称性的基础表示中。暂时略去自旋和味量子数，我们可以用 $q^i$ 表示夸克，其中 $i = 1, 2, 3$ 是色指标。于是夸克在色 $SU(3)$ 对称性的基础表示或「3」表示下变换。反夸克 $q_i$ 在 $\bar{3}$ 表示下变换。一个 3 与一个 $\bar{3}$ 的内积是 $SU(3)$ 的不变量。还可以用三个 3 的完全反对称组合 $\epsilon^{ijk}$ 构成不变量：该对象在幺正变换下按
\begin{equation}
  \epsilon^{i'j'k'}U_{ii'}U_{jj'}U_{kk'} = (\det U)\epsilon^{ijk} ,
  \label{eq:17.2}
\end{equation}
变换，因此在 $\det U = 1$ 的 $SU(3)$ 变换下不变。在强子波函数在 $SU(3)$ 对称变换下必须不变的假定下，这两种组合是仅有的简单组合：
\begin{equation}
  \bar{q}_i q^i,\qquad \epsilon_{ijk}q^i q^j q^k,\qquad
  \epsilon^{ijk}\bar{q}_i\bar{q}_j\bar{q}_k .
  \label{eq:17.3}
\end{equation}
也就是说，物理强子在色下是单态的假设意味着唯一可能的轻强子是介子、重子和反重子。

与原始夸克模型一样，色假设在唯象上取得了成功，但也提出了额外的问题：为什么夸克应当具有这个看似多余的属性，又是什么机制确保所有强子波函数都是色单态？这些问题的答案并非来自强子谱学，而是来自第~14 章描述的深度非弹性散射实验，以及随后对具有渐近自由性质的部分子束缚理论的探索。当发现非阿贝尔规范理论具有这一性质时，剩下要做的就是确定正确的规范群和费米子表示。由于色对称没有其他明显的物理作用，自然就把这种对称性与规范群等同起来，色就是夸克的规范量子数。这种推理导致了一个强相互作用模型，其中强相互作用是各种味的夸克系统，每个夸克都被指定到定域规范群 $SU(3)$ 的基础表示。$SU(3)$ 规范场的量子称为\emph{胶子}，这个理论被称为量子色动力学，即 QCD。

在本书中，我们主要研究 QCD 在耦合常数已经变小的能区中的性质。然而，我们应该指出，也可以利用 Wilson 引入的近似方案在强耦合区研究 QCD，在该方案中，连续规范理论被四维 Euclidean 格点上的离散统计力学系统所取代。利用这一近似，Wilson 证明了，对于足够强的耦合，QCD 表现出色的\emph{禁闭}：理论中仅有的有限能量渐近态是色 $SU(3)$ 的单态。于是，那个解释强子谱的特设假设结果是非阿贝尔规范理论与色耦合的一个推论。

\begin{figure}[t]
  \centering
  \includegraphics[width=0.5\textwidth]{fig17.1.png}
  \caption{强耦合规范理论中与色源分离相联系的规范电场位形。}
  \label{fig:17.1}
\end{figure}

如果试图把一个色单态分离成带色的组分——例如，把一个介子离解成夸克和反夸克——那么两个源之间就会形成一条规范场管，如图~\ref{fig:17.1} 所示。在耦合足够强的非阿贝尔规范理论中，这条管具有固定的半径和能量密度，因此分离色源的能量代价与分离距离成正比地增长。这种类型的力定律可以在短距离处弱而在长距离处强，从而解释了观测不到孤立夸克这一事实。我们将在尾声部分进一步讨论规范理论的大距离、强耦合极限。

量子色动力学的短距离极限可以用我们在前面各章发展的费曼图技术来方便地研究。这里渐近自由使耦合变弱，存在一个从自由夸克和胶子模型出发的合理的图解微扰论。以下几节讨论可以在高能实验中观测到的夸克与胶子之间的基本相互作用。

\section{$e^{+}e^{-}$ 湮灭为强子}

涉及夸克的最简单反应是 $e^{+}e^{-}$ 湮灭中夸克对的产生，这是我们在第~5.1 节已经处理过的过程。在那里我们只在最基础的层次上分析这一过程，把它看成自由夸克由虚光子产生的纯 QED 反应。这一最低阶过程的图如图~\ref{fig:17.2}(a) 所示。总截面的计算包括对夸克场各种色态的求和，因此证实了允许的色数为 3。把色因子与夸克电荷平方结合起来，我们得到（式~\eqref{eq:5.16}）
\begin{equation}
  \sigma(e^{+}e^{-}\to\text{强子}) = \sigma_0\cdot 3\cdot
    \sum_f Q_f^2 ,
  \label{eq:17.4}
\end{equation}
其中 $\sigma_0$ 是 $e^{+}e^{-}\to\mu^{+}\mu^{-}$ 的 QED 截面，
\begin{equation}
  \sigma_0 = \frac{4\pi\alpha^2}{3s} ,
  \label{eq:17.5}
\end{equation}
而 \eqref{eq:17.4} 中的求和遍及所有夸克味。这个公式假定质心能量足够高，以致我们可以忽略夸克质量。

\begin{figure}[t]
  \centering
  \includegraphics[width=0.5\textwidth]{fig17.2.png}
  \caption{对 QCD 中 $e^{+}e^{-}\to$ 强子过程有贡献的图：(a) 领头阶图；(b) $\alpha_s$ 阶修正。}
  \label{fig:17.2}
\end{figure}

当我们把夸克耦合到 $SU(3)$ 规范理论时，我们增加了许多重要过程，它们既影响这个截面的值，也影响它所包含的末态。一些最重要的效应无法在微扰论的框架内讨论。特别是，尽管领头图含有自由夸克，反应中涌现出的粒子却是色单态介子和重子。然而，我们将发现，含有夸克和胶子的 QCD 微扰论确实对 $e^{+}e^{-}$ 湮灭为强子作出一些重要预言。我们在得到这些预言的过程中发展的思想也将适用于许多其他强相互作用过程。

\subsection{总截面}

由胶子交换和发射引起的 $e^{+}e^{-}$ 湮灭率的领头修正如图~\ref{fig:17.2}(b) 所示。这些正是第~I 部分末项工程中计算的图。前两个图给出除夸克和反夸克之外再产生一个胶子的、量级为 $g^2$ 的截面，其中 $g$ 是 $SU(3)$ 规范耦合。第三个图必须在振幅中与领头图相加，以给出无胶子发射时 $q\bar{q}$ 产生率的修正。在第~I 部分中，我们就像强相互作用是阿贝尔规范理论那样计算了这两项贡献。为得到 QCD 中相应的表达式，我们只需把阿贝尔公式乘以群论因子
\begin{equation}
  \mathrm{tr}[t^at^a] = C_2(r)\cdot\mathrm{tr}[\mathbb{1}]
  = \frac{4}{3}\cdot 3 ,
  \label{eq:17.6}
\end{equation}
其中我们用式~(15.97) 计算了 $SU(3)$ 基础表示的 $C_2(r)$。因子 3 就是式~\eqref{eq:17.4} 中出现的同一个色求和。于是我们可以通过作替换
\begin{equation}
  g^2 \to \frac{4}{3}g^2,\qquad \text{或}\qquad
  \alpha_g \to \frac{4}{3}\alpha_s ,
  \label{eq:17.7}
\end{equation}
从末项工程的公式得到 QCD 的正确公式，其中
\begin{equation}
  \alpha_s = \frac{g^2}{4\pi}
  \label{eq:17.8}
\end{equation}
是精细结构常数的强相互作用类比量。

第~I 部分末项工程的最终结果是 $e^{+}e^{-}$ 湮灭中产生强子的总截面公式。如果我们把 $\alpha_g$ 换成 $(4/3)\alpha_s$，该结果变为
\begin{equation}
  \sigma(e^{+}e^{-}\to\text{强子}) = \sigma_0\cdot 3\cdot
    \sum_f Q_f^2\ \Bigl[ 1 + \frac{\alpha_s}{\pi} + \mathcal{O}(\alpha_s^2)
      \Bigr] .
  \label{eq:17.9}
\end{equation}
这实际上是两个基本过程——$e^{+}e^{-}\to q\bar{q}$（包括图~\ref{fig:17.2}(b) 中第三个图的修正）和 $e^{+}e^{-}\to q\bar{q}g$——的速率之和。尽管当胶子质量趋于零时这些过程中每一个的速率都是发散的，但当它们合并在一起时该发散相互抵消。这是我们在第~6.4 节和第~6.5 节以电子散射为例研究过的红外发散抵消现象的又一个例子。在那里我们证明了末态通过软光子和共线光子发射的修饰并不影响总的散射率。这里我们再次看到，红外发散在总速率中抵消，尽管对实胶子和虚胶子修正的求和留下一个简单的数值修正。

直观地理解红外对数项的抵消并不困难。原始过程 $e^{+}e^{-}\to q\bar{q}$ 极其迅速：由于虚光子离壳量为 $q^2 = s$，夸克在时间 $1/\sqrt{s}$ 内被产生。然而，共线胶子的发射，以及与软胶子交换相联系的虚修正，则在长得多的时标上发生。在含有胶子发射的图中，虚夸克或虚反夸克离质量壳的量为 $p^2_{\perp g}$，其中 $p_{\perp g}$ 是胶子相对于 $q\bar{q}$ 系统的横向动量。于是这个虚态在衰变之前存活时间 $1/p_{\perp g}$。这样慢的过程不可能影响 $q\bar{q}$ 对是否产生的概率；它只能影响 $q\bar{q}$ 系统将演化成的末态的性质。按照这种逻辑，唯一能影响总截面的微扰修正就是那些 $p_{\perp g} \sim \sqrt{s}$ 的修正。表达这一结论的另一种方式是论证，在对红外敏感区域的贡献抵消之后，剩下的贡献来自大实胶子或虚胶子动量区域。无论用哪种论证，公式 \eqref{eq:17.9} 都应该是 QCD 微扰论的一个有意义的预言，即使它涉及对软胶子发射区域的积分。

\subsection{$\alpha_s$ 的跑动}

公式 \eqref{eq:17.9} 涉及重整化耦合 $\alpha_s$，它依赖于重整化标度 $M$。为得到独立于这一任意标度的预言，我们使用重整化群。QCD 的 $\beta$ 函数在第~16 章中计算过：
\begin{equation}
  \beta(g) = -g\,\frac{b_0}{16\pi^2}\,g^2 ,
  \label{eq:17.11}
\end{equation}
其中
\begin{equation}
  b_0 = 11 - \frac{2}{3}n_f ,
  \label{eq:17.12}
\end{equation}
适用于具有 $n_f$ 味夸克的 $SU(3)$。重整化群方程
\begin{equation}
  \frac{d\alpha_s}{d\log(Q/M)} = \beta(\alpha_s)
\end{equation}
的解是
\begin{equation}
  \alpha_s(Q) = \frac{\alpha_s(M)}{1 + (b_0\alpha_s/2\pi)\log(Q/M)} .
  \label{eq:17.14}
\end{equation}
$\sigma$ 对 $\alpha_s$ 的显式依赖可以通过把 $f(\alpha_s(\sqrt{s}))$ 展开中的逐项与微扰展开中的各项相匹配来找到。到第一修正的阶，我们简单地得到
\begin{equation}
  \sigma = \sigma_0\cdot 3\cdot \sum_f Q_f^2\ \Bigl[ 1 + \frac{\alpha_s
    (\sqrt{s})}{\pi} + \cdots \Bigr] .
  \label{eq:17.15}
\end{equation}
于是 Callan-Symanzik 方程指示我们把固定的重整化耦合 $\alpha_s$ 换成在 $Q^2 = s$ 处求值的跑动耦合常数 $\alpha_s(Q)$。

由于固定耦合 $\alpha_s$ 依赖于任意重整化点 $M$，有时把从公式中完全消去它是有用的。为此，我们定义一个质量标度，按惯例称之为 $\Lambda$（不要与紫外截断混淆！），它满足
\begin{equation}
  1 = \bar{g}^2(b_0/8\pi^2)\log(M/\Lambda) .
  \label{eq:17.16}
\end{equation}
于是式~\eqref{eq:17.14} 可以重排成
\begin{equation}
  \alpha_s(Q) = \frac{4\pi}{b_0\log(Q^2/\Lambda^2)} .
  \label{eq:17.17}
\end{equation}
这个公式最清楚地表达了如下陈述：对大 $Q$，$\alpha_s(Q)$ 按 $(\log Q)^{-1}$ 变小。动量标度 $\Lambda$ 是当 $Q^2$ 减小时 $\alpha_s$ 变强的标度。

对这个反应及其他反应速率的实验测量给出 $\Lambda \approx 200$ MeV。QCD 微扰论只在 $Q$ 比它稍大时才有效，比如在 $Q = 1$ GeV 以上，那里 $\alpha_s(Q) \sim 0.4$。强相互作用在大于 $\sim 1/\Lambda$ 的距离处变强，这大致是轻强子的大小。

尽管 $e^{+}e^{-}$ 湮灭截面的例子特别简单，因为它只依赖一个动量不变量，类似的结论也适用于其他 QCD 预言。在分析对夸克和胶子亚结构敏感的强相互作用过程时，我们将找到依赖重整化耦合 $\alpha_s$ 的反应截面领头阶公式。为使这些表达式满足 Callan-Symanzik 方程，我们必须把这一固定耦合换成跑动耦合常数 $\alpha_s(Q)$，其 $Q$ 取反应动量不变量的量级。由于跑动耦合常数只对数地依赖 $Q$，我们不必担心精确选择 $Q$。如果我们把正确的标度猜错一个因子 2，这在 $\alpha_s(Q)$ 中引起的误差量级为 $(\log Q)^{-2} \sim \alpha_s^2(Q)$。反过来，这个歧义可以通过计算到 $\alpha_s$ 的下一阶来消除。

在结束对 $e^{+}e^{-}$ 湮灭截面的这种形式处理之前，我们应该补充一个保留条件。在第~12.2 节的开头，我们曾指出重整化群预言可能因物理阈及其相关奇异性的出现而变得复杂，因此我们只在相关动量不变量 $P^2$ 大且类空时才陈述这些预言。在本章中，我们将关注夸克和胶子反应的截面，它们在壳上求值。这引入了额外的原则性复杂性。例如，为了把 Callan-Symanzik 方程应用于 $\sigma(s)$，我们需要知道这个量不含红外发散，否则其正规化子可能提供另一个有量纲参数。在整个本章中，我们将假设与软胶子和共线胶子相联系的类似发散抵消发生在我们所关注的过程中。QCD 中这些抵消的完整证明是可以完成的，但相当技术化。%
\footnote{关于证明微扰 QCD 公式合理性的定理综述，参见 J.~C.~Collins 和 D.~E.~Soper，\emph{Ann.\ Rev.\ Nucl.\ Sci.} \textbf{37}, 383 (1987)。} 在某些情形下，另一种观点是可能的；人们可以通过把在壳振幅与在类空区域求值的 Green 函数相联系，来证明用重整化群分析它的合理性。这种分析方法带来了自身的洞见，将是第~18 章的主要论题。

\subsection{胶子发射与喷注产生}

第~I 部分末项工程的第二个结果是带胶子发射的 $q\bar{q}$ 产生的微分截面公式。用 \eqref{eq:17.7} 把这个公式转写到 QCD 给出如下结果：设 $x_1$、$x_2$、$x_3$ 分别为夸克、反夸克和胶子能量与电子束能量的比值。它们满足 $0 < x_i < 1$ 和 $x_1 + x_2 + x_3 = 2$。于是 $e^{+}e^{-}\to q\bar{q}g$ 的截面由
\begin{equation}
  \frac{d^2\sigma}{dx_1\,dx_2}(e^{+}e^{-}\to q\bar{q}g)
  = \sigma_0\cdot\Bigl( 3\sum_f Q_f^2 \Bigr)\cdot
    \frac{2\alpha_s}{3\pi}\ \frac{x_1^2+x_2^2}{(1-x_1)(1-x_2)} .
  \label{eq:17.18}
\end{equation}
给出。当 $x_1$ 或 $x_2$ 趋于 1 时，这个截面是奇异的。极限 $x_1\to 1$ 对应于夸克具有最大可能能量、而反夸克和胶子沿相反方向飞出并分享剩余能量的位形。于是反夸克和胶子具有几乎共线的类光动量矢量，从而形成一个不变质量非常小的系统。类似地，极限 $x_2\to 1$ 对应于夸克和胶子共线的位形。这些奇异性是积分截面在胶子质量趋于零的极限下发散的原因。

我们应该如何解释这些奇异性？在第~6.1 节对轫致辐射的一般处理中，我们看到散射电子发射光子对于共线辐射增强一个量级为 $\log(s/m^2)$ 的因子，其中 $m$ 是电子质量。于是发射共线光子的总速率在零质量极限下形式上是发散的。同样的结论也适用于夸克发射胶子。在零质量极限下共线发射出现的发散称为\emph{质量奇异性}。在 QED 中，我们看到当 $q^2$ 大时，质量奇异性标志着强共线辐射的一个真实物理效应。在 QCD 中，我们可以预期在这个极限下有强胶子辐射，但我们必须仔细思考这种辐射在实验中如何显现。无论共线胶子是否被辐射，从反应中涌现出的夸克和反夸克都将与其他产物发生进一步的软相互作用。这些过程必须继续下去，产生夸克-反夸克对并发射和吸收胶子，直到所有带色粒子都被收集进色单态强子。因此，一个或多个共线胶子的存在对末态没有显著影响，末态由两个背对背的强子喷注组成。因此，当胶子横向动量小于软胶子相互作用的典型标度（约为 1 GeV）时，公式 \eqref{eq:17.18} 没有用处。

然而，当胶子以相对于 $q\bar{q}$ 轴显著的横向动量被发射时，随后的软交换不可能收回或反转这个横向动量。在这种情况下，$q\bar{q}g$ 系统演化成三个不同强子喷注的系统。因此，在离共线区域足够远处，我们可以把式~\eqref{eq:17.18} 解释为产生具有三个强子喷注事件——其能量为电子束能量的 $x_1$、$x_2$、$x_3$ 倍——的截面。

通过与上面为总截面给出的分析类似的分析，我们可以把固定耦合常数 $\alpha_s$ 换成跑动的 $\alpha_s(Q)$ 来改进式~\eqref{eq:17.18}。对 $Q$ 的一个合理选择是式~\eqref{eq:17.9} 下方描述的胶子横向动量 $p_{\perp g}$。然而，如果这个横向动量太小，$\alpha_s(Q)$ 就会很大，我们的领头阶公式将失效。这是我们不能在胶子横向动量小于约 1 GeV 时使用公式 \eqref{eq:17.18} 的第二个原因。

另一方面，当胶子横向动量远大于 1 GeV 时，没有理由怀疑 QCD 微扰论。软过程不可能干扰强子态的三喷注性质，渐近自由保证了耦合常数很小，因此微扰论的领头阶将是好的近似。

三喷注截面 \eqref{eq:17.18} 是从 QCD 中使用微扰论得到的那类预言的一个好例子。我们用强相互作用过程传递给强子组分的动量转移 $Q$ 来描述该过程。QCD 微扰论对这类反应中能量和动量流入末态强子喷注系统的流动作出预言。如果 $Q$ 小，微扰论失效，我们得不到有用的预言。然而，如果 $Q$ 大，QCD 的渐近自由意味着夸克和胶子的费曼图将正确地预言末态强子喷注系统的行为。

\section{深度非弹性散射}

在 $e^{+}e^{-}$ 湮灭为强子之后，涉及强相互作用粒子的次简单反应是电子对质子或其他强子的散射。在最基础的层次上，这个反应可以被看成电子对质子内部夸克的电磁散射。%
\footnote{电子换成缪子也同样可以；所有相同公式都适用于这种情况。轻子也可以经由中性流弱相互作用散射夸克，我们将在第~20 章看到这一点。} 图~\ref{fig:17.3} 给出了一种把这个过程形象化的方式。称质子动量为 $P$，初始夸克动量为 $p$。称电子初动量和末动量分别为 $k$ 和 $k'$。如果测量末态电子动量，就可以推知虚光子传递给强子系统的动量 $q = k-k'$。矢量 $q$ 是类空的，人们按惯例写
\begin{equation}
  q^2 = -Q^2 .
\end{equation}
如果不变量动量转移 $Q^2$ 大，夸克以随后的软过程无法平衡的方式从质子中射出。然而，这些软过程将产生胶子和夸克-反夸克对，它们最终中和色，并使被撞击的夸克在来自电子的动量转移方向上以强子喷注的形式物质化。通常末态强子系统的总不变质量很大，因为被撞击的夸克相对于其他「旁观者」夸克携带很大的动量。在这种情况下，该过程被称为\emph{深度非弹性散射}。

\begin{figure}[t]
  \centering
  \includegraphics[width=0.5\textwidth]{fig17.3.png}
  \caption{QCD 中的深度非弹性散射。图中显示了一个高能电子散射取自质子波函数的夸克时的动量流。}
  \label{fig:17.3}
\end{figure}

为推导电子-质子散射截面的第一近似，我们从一个电子和质子快速相向运动的参考系来考虑这个反应，例如电子-质子质心系。我们假定质心能量足够大，以致在处理运动学时可以忽略质子质量。于是质子沿碰撞轴具有几乎类光的动量。质子的组分也具有类光动量，它们与质子的动量几乎共线。这是因为组分只有通过硬胶子交换才能获得大的横向动量，而这一过程由于 $\alpha_s$ 在大动量标度下很小而受到抑制。于是，在 QCD 微扰论的领头阶，我们可以写
\begin{equation}
  p = \xi P ,
  \label{eq:17.19}
\end{equation}
其中 $\xi$ 是 0 到 1 之间的数，称为组分的\emph{纵分数}。在 $\alpha_s$ 的领头阶，我们还可以忽略碰撞过程中的胶子发射或交换。于是电子-质子散射的截面由给定纵分数 $\xi$ 处的电子-夸克散射截面乘以质子在该 $\xi$ 值处含有夸克的概率、再对 $\xi$ 积分给出。

但是，质子含有某个具有某个动量分数的组分的概率无法用 QCD 微扰论计算，因为它依赖于决定质子作为夸克和胶子束缚态结构的软过程。因此我们将把这一概率视为未知函数，由实验来确定。最终，我们将需要对质子波函数中可能找到的每一种夸克、反夸克和胶子都使用这样一个概率函数。这些组分统称为\emph{部分子}。对每种部分子 $f$，我们把在质子中找到纵分数为 $\xi$ 的 $f$ 类组分的概率写成
\begin{equation}
  f_f(\xi)\,d\xi = \text{找到一个纵分数为 } \xi \text{ 的 } f \text{ 类部分子的概率} .
  \label{eq:17.20}
\end{equation}
函数 $f_f(\xi)$ 称为\emph{部分子分布函数}。使用这种记号，电子-质子非弹性散射的截面在 $\alpha_s$ 的领头阶由表达式
\begin{equation}
  \sigma\bigl( e^{-}(k)p(P) \to e^{-}(k') + X \bigr)
  = \int_0^1 d\xi\ \sum_f f_f(\xi)\ \hat{\sigma}\bigl( e^{-}(k)q_f(\xi P)
    \to e^{-}(k') + q_f(p') \bigr) ,
  \label{eq:17.21}
\end{equation}
给出，其中 $X$ 代表任意强子末态。\eqref{eq:17.21} 中的求和包含来自组分反夸克以及组分夸克的贡献。

方程 \eqref{eq:17.21} 等价于我们在第~14 章为这个反应构造的公式 \eqref{eq:14.8}。现在我们看到，这个公式的合理性由 QCD 耦合常数在大动量标度下很小来保证。然而，重要的是记住，\eqref{eq:17.21} 不是 QCD 的完整预言，而只是 $\alpha_s$ 展开的第一项；这一近似层次称为\emph{部分子模型}。对式~\eqref{eq:17.21} 的高阶 QCD 修正既涉及对电子-夸克散射截面的修改，也涉及对部分子分布函数的修改。这些修正中最重要的将在第~17.5 节讨论。

以同样的方式，质子涉及大动量转移的所有其他反应也都具有部分子模型描述。在 QCD 中，所有这些反应截面都由夸克和胶子的散射振幅计算。任何过程中部分子的初始运动都由深度非弹性散射中出现的相同部分子分布函数 $f_f(\xi)$ 描述。

现在让我们复习第~14 章的分析，得出深度非弹性散射截面的显式领头阶公式。在式~\eqref{eq:14.3} 中，我们写下了部分子层次过程的领头阶微分截面，
\begin{equation}
  \frac{d\sigma}{dt}\bigl( e^{-}q \to e^{-}q \bigr)
  = \frac{2\pi\alpha^2 Q_f^2}{s^2}\ \frac{s^2 + u^2}{t^2} .
  \label{eq:17.22}
\end{equation}
一般地，我们将用符号 $s$、$t$、$u$ 表示部分子层次二体散射过程的 Mandelstam 变量。这些变量必须与强子系统或散射电子的可观测性质相联系。对无质量的初末态粒子，
\begin{equation}
  s + t + u = 0 .
\end{equation}
在深度非弹性散射的情形下，
\begin{equation}
  t = -Q^2
  \qquad\text{和}\qquad
  s = 2p\cdot k = 2\xi P\cdot k = \xi s .
\end{equation}
把这些关系与 \eqref{eq:17.22} 以及部分子模型因子化 \eqref{eq:17.21} 结合起来，我们得到深度非弹性散射截面：
\begin{equation}
  \frac{d^2\sigma}{dx\,dQ^2}\bigl( e^{-}p \to e^{-}X \bigr)
  = \frac{2\pi\alpha^2}{Q^4}\ \Bigl[ 1 + (1-y)^2 \Bigr]\ \sum_f Q_f^2\,
    f_f(x) ,
  \label{eq:17.30}
\end{equation}
其中 $x = Q^2/2P\cdot q$ 是 Bjorken 标度变量，$y = 1 - Q^2/xs$。这个公式展示了 Bjorken 标度性质，即结构函数只依赖 $x$ 而不依赖 $Q^2$，直至 $\alpha_s$ 阶的修正。表达式 \eqref{eq:17.30} 还预言了对 $y$ 的特定依赖：因子 $[1 + (1-y)^2]$ 是电子对无质量费米子散射的特征。在部分子与夸克之间的关系仍不清楚的时候，这个关系为深度非弹性散射中所涉及的部分子是费米子提供了证据。

\subsection{中微子深度非弹性散射}

由于式~\eqref{eq:17.30} 中对夸克味的求和因子化，人们无法仅从电子散射实验确定各个部分子分布函数 $f_f(x)$。然而，人们可以通过中微子深度非弹性散射获得关于质子结构的更详细信息。

中微子电荷为零，因此不通过光子交换相互作用，但它们通过弱相互作用与夸克相互作用。我们将在第~20 章详细讨论弱相互作用；目前，让我们采用一个集中于图~\ref{fig:17.4} 所示基本过程的简化描述。在这个过程中，中微子转化为相应的带电轻子，%
\footnote{还存在一种中性流弱相互作用，其中中微子仍然是中微子；参见习题~20.4。} 同时交换一个虚的大质量矢量玻色子 $W^{+}$。这个玻色子与把 $d$ 夸克转化为 $u$ 夸克的夸克流耦合。这一交换过程的效果是提供一种不同的、但被完全表征的注入硬动量转移 $q$ 的方法。这一过程的振幅由有效拉格朗日量
\begin{equation}
  \Delta\mathcal{L} = -\frac{g^2}{m_W^2}\ \bar{\ell}\gamma^\mu
    \frac{1-\gamma^5}{2}\nu\ \bar{u}\gamma_\mu\frac{1-\gamma^5}{2}d
    + \mathrm{h.c.},
  \label{eq:17.31}
\end{equation}
描述，其中 $\ell$、$\nu$、$d$、$u$ 是与带电轻子、中微子以及 $d$ 和 $u$ 夸克相联系的费米子场，$g$ 是弱相互作用耦合常数。因子 $1/m_W^2$ 来自 $W$ 玻色子传播子，在 $q^2 \ll m_W^2$ 的极限下考虑。前两个因子常常用 Fermi 常数 $G_F$ 来写，定义为
\begin{equation}
  \frac{G_F}{\sqrt{2}} = \frac{g^2}{8m_W^2} .
  \label{eq:17.32}
\end{equation}
这个常数给出远小于 $m_W$ 的能量下弱相互作用的强度。弱相互作用的关键性质——在 \eqref{eq:17.31} 中显式展示——是 $W$ 玻色子只与相对论性费米子的\emph{左手}螺旋态耦合。这一性质的更深层意义将在第~20 章讨论。

\begin{figure}[t]
  \centering
  \includegraphics[width=0.5\textwidth]{fig17.4.png}
  \caption{由弱相互作用中介的基本中微子散射过程。}
  \label{fig:17.4}
\end{figure}

出于技术原因，用缪子中微子做中微子深度非弹性散射最容易，缪子中微子在发射 $W$ 玻色子后转化为缪子。让缪子反中微子散射核靶同样是可行的，而且，正如我们将看到的，比较中微子和反中微子的效应是很有趣的。由于质子含有少量较重夸克 $s$、$c$ 的掺混，它们也对中微子深度非弹性散射给出小的贡献。然而，在我们的讨论中我们将忽略这些贡献。

中微子深度非弹性散射的截面由与 \eqref{eq:17.21} 类似的公式给出，只是把光子交换截面换成由 $W$ 交换产生的截面。直接算出这个截面是直截了当的。然而，如果我们回顾第~5 章并回想这个方程的结构如何由各种螺旋度贡献产生，我们也可以从式~\eqref{eq:17.22} 得到该结果。在 \eqref{eq:17.22} 中，分母中的因子 $t^2$ 来自光子传播子；在弱相互作用情形下，这个因子被 $m_W^4$ 替换。因子 $\tfrac{1}{2}s^2$ 来自 Dirac 矩阵代数。我们在第~5.2 节看到，第一项是左手电子散射左手费米子或右手电子散射右手费米子的贡献，第二项来自其他螺旋度组合。对中微子-夸克散射的情形，相互作用 \eqref{eq:17.31} 只允许左手中微子散射左手夸克，因此只出现 $s^2$ 项。为确定截面的整体归一化，我们注意到，由于中微子在弱相互作用中产生，它们总是具有左手极化，因此不应做极化平均。另一方面，我们仍然必须对初始夸克的极化求平均。总的说来，我们得到
\begin{equation}
  \frac{d\sigma}{dt}\bigl( \nu_\mu d \to \mu^{-}u \bigr)
  = \frac{G_F^2}{2(4\pi)^2}\ \frac{s^2}{t^2}\ \cdots
  = \frac{G_F^2}{2\pi}\,s\,(1-y)^2 .
  \label{eq:17.33}
\end{equation}
从 \eqref{eq:17.31} 出发作显式计算来检验这个公式是容易的。

反中微子散射夸克的截面可以用同样的方法算出。注意这个反应涉及 $W^{-}$ 的交换，因此把 $u$ 夸克转化为 $d$ 夸克。然而，$u$ 夸克仍然必须是左手的。与上一段的唯一不同在于如下事实：与相互作用 \eqref{eq:17.31} 耦合的反中微子是右手的，因此截面来自 \eqref{eq:17.22} 中正比于 $u^2$ 的项：
\begin{equation}
  \frac{d\sigma}{dt}\bigl( \bar{\nu}_\mu u \to \mu^{+}d \bigr)
  = \frac{G_F^2}{2(4\pi)^2}\ \frac{u^2}{s^2}\ \cdots
  = \frac{G_F^2}{2\pi}\,s\,(1-y)^2 .
  \label{eq:17.34}
\end{equation}
同样，直接验证这个公式是容易的。中微子-反夸克散射（把 $\bar{u}$ 转化为 $\bar{d}$）的截面也由式~\eqref{eq:17.34} 给出，而反中微子-反夸克散射（把 $\bar{d}$ 转化为 $\bar{u}$）的截面由式~\eqref{eq:17.33} 给出。

为把这些部分子层次截面转换为物理截面，我们把它们与部分子分布函数结合起来。运动学与电子散射的情形完全相同。于是，遵循导出式~\eqref{eq:17.30} 的论证，我们得到表达式
\begin{equation}
  \begin{aligned}
    \frac{d^2\sigma}{dx\,dy}\bigl( \nu p \to \mu^{-}X \bigr)
    &= \frac{G_F^2}{\pi}\,s\ \bigl[ xf_d(x) + xf_u(x)\cdot(1-y)^2 \bigr],\\
    \frac{d^2\sigma}{dx\,dy}\bigl( \bar{\nu} p \to \mu^{+}X \bigr)
    &= \frac{G_F^2}{\pi}\,s\ \bigl[ xf_u(x)\cdot(1-y)^2 + xf_d(x) \bigr].
  \end{aligned}
  \label{eq:17.35}
\end{equation}
根据这些关系，中微子深度非弹性散射使人们能够分别描绘质子中 $u$ 和 $d$ 夸克与反夸克的部分子分布函数。

此外，式~\eqref{eq:17.35} 作出一个引人注目的定性预言：只要质子（或中子）由夸克构成且额外夸克-反夸克对很少，中微子深度非弹性散射截面在 $y$ 上应该是常数，而反中微子散射截面应该按 $(1-y)^2$ 下降。这些深度非弹性截面测量到的 $y$ 依赖如图~\ref{fig:17.5} 所示。部分子描述预言的定性行为清晰可见；与严格预言的偏离可以用核子波函数中的小反夸克组分来解释。

\begin{figure}[t]
  \centering
  \includegraphics[width=0.6\textwidth]{fig17.5.png}
  \caption{由 CDHS 实验（J.~G.~H.~de~Groot 等人，\emph{Z.\ Phys.\ C} \textbf{1}, 143 (1979)）测量的中微子和反中微子对铁靶深度非弹性散射的 $y$ 分布。实线是对 $A + B(1-y)^2$ 形式的拟合。}
  \label{fig:17.5}
\end{figure}

\subsection{部分子分布函数}

既然电子和中微子深度非弹性散射的部分子模型预言确实拟合了数据，人们就可以利用这些关系提取部分子分布函数，从而了解质子结构的一些信息。%
\footnote{关于从数据提取部分子分布函数的详细讨论，参见 G.~Sterman 等人，\emph{Rev.\ Mod.\ Phys.} \textbf{67}, 157 (1995)。} 一组经过选择以拟合所有可用数据的分布函数如图~\ref{fig:17.6} 所示。由于所有这些分布，特别是反夸克的分布，在小 $x$ 处尖锐地达到峰值，我们对每种都绘制了 $xf_f(x)$。正如我们在第~14 章指出的，实验中观测到 Bjorken 标度的小破坏，因此这些分布函数随 $Q^2$ 缓慢变化。图中显示了这些函数在 $Q^2 = 4\ \mathrm{GeV}^2$ 处的值。我们将在第~17.5 节看到，这种 Bjorken 标度破坏是 QCD 高阶修正的效应；我们还将论证，对这种标度破坏的测量使人们能够确定胶子的部分子分布函数 $f_g(x)$。预期到这一结果，我们也在图中绘制了这个函数。毫不奇怪，人们发现 $u$ 和 $d$ 夸克最有可能携带质子的很大一部分动量，而反夸克和胶子倾向于具有小的纵分数。

\begin{figure}[t]
  \centering
  \includegraphics[width=0.6\textwidth]{fig17.6.png}
  \caption{质子中夸克、反夸克和胶子在 $Q^2 = 4\ \mathrm{GeV}^2$ 处的部分子分布函数 $xf_f(x)$。这些分布来自 CTEQ 合作组（CTEQ2L）对深度非弹性散射数据所作的拟合，见 J.~Botts 等人，\emph{Phys.\ Lett.\ B} \textbf{304}, 159 (1993)。}
  \label{fig:17.6}
\end{figure}

由于部分子分布是在质子中找到各种组分的概率，它们必须以反映质子量子数的方式归一化。质子是 $uud$ 的束缚态，加上一些夸克-反夸克对的掺混。因此它应比相应的反夸克多出两个 $u$ 夸克和一个 $d$ 夸克。这些考虑意味着约束
\begin{equation}
  \int_0^1 dx\ \bigl[ f_u(x) - f_{\bar{u}}(x) \bigr] = 2,\qquad
  \int_0^1 dx\ \bigl[ f_d(x) - f_{\bar{d}}(x) \bigr] = 1 .
  \label{eq:17.36}
\end{equation}

到目前为止，我们只讨论了质子的部分子分布。然而，类似的考虑适用于任何其他强子。每个强子都有自己的部分子分布函数集合；它们服从与 \eqref{eq:17.36} 类似的求和规则，但反映强子的特定量子数。部分子分布函数还应该反映联系不同强子的对称性。例如，由于中子可以通过交换质子中 $u$ 和 $d$ 夸克的角色产生，精确到百分之几，它的分布函数服从
\begin{equation}
  f_u^n(x) = f_d(x),\qquad f_d^n(x) = f_u(x),\qquad
  f_s^n(x) = f_s(x),\qquad \text{等等}.
  \label{eq:17.37}
\end{equation}
在这些方程中，以及此后，没有特殊标号的分布函数都指质子。反质子的部分子分布函数由精确关系
\begin{equation}
  f_u^{\bar{p}}(x) = f_{\bar{u}}(x),\qquad
  f_{\bar{u}}^{\bar{p}}(x) = f_u(x),\qquad \text{等等}.
  \label{eq:17.38}
\end{equation}
给出。

无论如何，部分子携带的总动量必须加起来等于强子的总动量。这意味着
\begin{equation}
  \int_0^1 dx\ x\ \bigl[ f_u(x) + f_d(x) + f_{\bar{u}}(x) + f_{\bar{d}}(x)
    + f_g(x) \bigr] = 1 .
  \label{eq:17.39}
\end{equation}
从深度非弹性散射数据提取的质子中夸克和反夸克的分布函数只贡献这个积分所需总值的约一半。剩余的能量-动量必须由胶子携带。

\section{强子对撞中的硬散射过程}

如果在非常高的能量下让强子与其他强子对撞，大多数碰撞将只涉及组分夸克和胶子的软相互作用。这类相互作用无法用微扰 QCD 处理，因为当动量转移小时 $\alpha_s$ 很大。然而，在某些碰撞中，两个夸克或胶子将交换垂直于碰撞轴的大动量 $p_\perp$。于是，与深度非弹性散射一样，基本相互作用相对于强子波函数的内部时标发生得非常迅速，因此最低阶 QCD 预言应该准确地描述该过程。同样，我们应该找到一个由领头阶子过程截面构成、并与部分子分布函数积分起来的部分子模型公式。对质子-质子散射的情形，这些函数将与轻子-质子深度非弹性散射中测量到的相同。

例如，如果硬的部分子层次过程涉及夸克-反夸克散射到末态 $Y$，领头阶 QCD 预言取如下形式
\begin{equation}
  \sigma\bigl( p(P_1) + p(P_2) \to Y + X \bigr)
  = \int_0^1 dx_1 \int_0^1 dx_2\ \sum_f f_f(x_1)f_{\bar{f}}(x_2)\cdot
    \hat{\sigma}\bigl( q_f(x_1P) + \bar{q}_{\bar{f}}(x_2P) \to Y \bigr) ,
  \label{eq:17.40}
\end{equation}
其中求和遍及所有种类的夸克和反夸克——$u$、$d$、$\bar{u}$、$\bar{d}$，\ldots。（这里同样，$X$ 代表任意强子末态。）相同公式在适当修改分布函数后适用于任何其他强子-强子碰撞。如果按某种不变度量，$q\bar{q}$ 反应中转移了大动量，这个公式将是好的第一近似。在本节中，我们将讨论这类过程的几个例子。

\subsection{轻子对产生}

强子-强子硬散射过程的第一个也是最简单的例子是轻子对产生，也称为 Drell-Yan 过程。在这个过程中，一个强子中的夸克与另一个强子中的反夸克湮灭为一个虚光子，虚光子物质化为一对轻子：
\begin{equation}
  q\bar{q} \to \gamma^{*} \to \mu^{+}\mu^{-} .
\end{equation}
这是我们在第~5 章研究过的反应 $e^{+}e^{-}\to\mu^{+}\mu^{-}$ 的强子类比。$q\bar{q}\to\mu^{+}\mu^{-}$ 的部分子层次截面通过把 $e^{+}e^{-}$ 公式中的电子电荷换成夸克电荷得到：
\begin{equation}
  \hat{\sigma}\bigl( q\bar{q} \to \mu^{+}\mu^{-} \bigr)
  = Q_f^2\,\frac{4\pi\alpha^2}{3\hat{s}} .
  \label{eq:17.41}
\end{equation}
把它与部分子模型公式 \eqref{eq:17.40} 结合起来，我们得到 Drell-Yan 截面。轻子对的不变质量为
\begin{equation}
  M^2 = x_1x_2\,s ,
\end{equation}
轻子对的纵动量决定差 $x_1 - x_2$。因此轻子对产生的截面为
\begin{equation}
  \frac{d\sigma}{dM^2}\bigl( pp \to \mu^{+}\mu^{-} + X \bigr)
  = \frac{4\pi\alpha^2}{9sM^2}\ \int dx_1\,dx_2\ \sum_f Q_f^2\,
    \bigl[ f_f(x_1)f_{\bar{f}}(x_2) + f_{\bar{f}}(x_1)f_f(x_2) \bigr]\,
    \delta(x_1x_2 - M^2/s) ,
  \label{eq:17.48}
\end{equation}
其中 $x_1$ 和 $x_2$ 由运动学的约束给出。值得注意的是，Drell-Yan 过程的截面逐点地由可从深度非弹性散射推出的信息决定。遗憾的是，\eqref{eq:17.48} 所蕴含的两个过程之间的关系会收到一个 $\alpha_s(M)$ 阶的修正，该修正数值上很大，而且为了对照实验检验这个预言，必须把它包括进去。

\subsection{对产生的一般运动学}

在推导式~\eqref{eq:17.48} 时，我们使用了部分子层次过程的总截面 \eqref{eq:17.41}，并对出射轻子的角分布求了积分。原则上，我们可以保留角度信息，推导一个三阶微分分布。这将是二体部分子层次反应可能作出的最完整预言。从更一般的观点出发推出这类反应的运动学将是有用的。

在一般情形下，来自质子 1 的 1 类部分子散射来自质子 2 的 2 类部分子，产生 3 类和 4 类部分子，动量转移平方为 $\hat{t}$。这一般过程如图~\ref{fig:17.8} 所示。在 Drell-Yan 过程中，部分子 3 和 4 是轻子。但这些部分子也可以是夸克或胶子，它们以强子喷注的形式物质化。我们假设所有部分子都可以当作无质量的。在部分子变量中，这一过程的截面为
\begin{equation}
  \frac{d\sigma}{dx_1\,dx_2\,d\hat{t}}\bigl( p_1p_2 \to 3 + 4 + X \bigr)
  = f_1(x_1)f_2(x_2)\ \frac{d\hat{\sigma}}{d\hat{t}}\bigl( 1 + 2 \to 3 + 4
    \bigr) .
  \label{eq:17.49}
\end{equation}

\begin{figure}[t]
  \centering
  \includegraphics[width=0.5\textwidth]{fig17.8.png}
  \caption{一个一般的二体部分子散射过程。}
  \label{fig:17.8}
\end{figure}

现在让我们把这个公式转换为末态的可观测参数。在 QCD 的领头阶，部分子 3 和 4 的横向动量必须等大反向，但它们的纵动量不受约束。我们将取末态的三个参数为部分子横向动量的共同大小 $p_\perp$ 以及末态部分子的纵向快度 $y_3$、$y_4$，它们由公式
\begin{equation}
  E_i = p_\perp\cosh y_i;\qquad
  p_{i\parallel} = p_\perp\sinh y_i .
  \label{eq:17.50}
\end{equation}
定义。纵向快度 $y_i$ 给出粒子 $i$ 从它纵动量为零的参考系出发的 boost。%
\footnote{在强子对撞的文献中，$y_i$ 通常简称为快度，并默认限定于纵向 boost。} 回顾第~3.3 节，快度在共线 boost 下简单地相加。横向动量在纵向 boost 下不变。因此，$(y_3, y_4, p_\perp)$ 是一组关于沿碰撞轴 boost 具有方便 Lorentz 变换性质的变量。我们现在将看到，这三个参数以直截了当的方式与底层变量 $x_1$、$x_2$、$\hat{t}$ 相联系。

考虑相撞部分子的质心系。在这个参考系中的总能量为 $\sqrt{\hat{s}}$。让我们用下标 $*$ 表示在这个参考系中测量的其他量，例如 $\theta_*$ 表示部分子散射角。于是
\begin{equation}
  p_{3\parallel*} = \frac{1}{2}\sqrt{\hat{s}}\cos\theta_*,\qquad
  p_{3\perp*} = \frac{1}{2}\sqrt{\hat{s}}\sin\theta_* ,
  \label{eq:17.51}
\end{equation}
而 $p_{4*}$ 的取向正好相反。这个参考系也是部分子 3 和 4 的质心系，因此
\begin{equation}
  y_{3*} = -y_{4*} = y_* .
  \label{eq:17.52}
\end{equation}
由于快度按平移变换，我们可以解出 $y_*$ 以及到达该参考系所必须 boost 的快度 $Y$：
\begin{equation}
  y_* = \frac{1}{2}(y_3 - y_4),\qquad
  Y = \frac{1}{2}(y_3 + y_4) .
  \label{eq:17.53}
\end{equation}
散射角 $\theta_*$ 通过把 \eqref{eq:17.51} 与关系式 $E_* = p_\perp\cosh y_*$ 结合，由 $y_*$ 确定：
\begin{equation}
  \frac{p_\perp}{\frac{1}{2}\sqrt{\hat{s}}} = \frac{1}{\sin\theta_*}
  = \cosh y_* .
  \label{eq:17.54}
\end{equation}
于是 Mandelstam 变量
\begin{equation}
  \hat{s} = \frac{4p_\perp^2}{\sin^2\theta_*},\qquad
  \hat{t} = -\frac{1}{2}\hat{s}(1-\cos\theta_*)
  \label{eq:17.55}
\end{equation}
可以表示为
\begin{equation}
  \hat{s} = 4p_\perp^2\cosh^2 y_*,\qquad
  \hat{t} = -2p_\perp^2\cosh y_*\,e^{-y_*} .
  \label{eq:17.56}
\end{equation}
我们可以把这两个表达式中的第一个与 $\hat{s} = x_1x_2s$ 结合来确定 $x_1$ 和 $x_2$：
\begin{equation}
  x_1 = \frac{2p_\perp}{\sqrt{s}}\,\cosh y_*\,e^{Y},\qquad
  x_2 = \frac{2p_\perp}{\sqrt{s}}\,\cosh y_*\,e^{-Y} .
  \label{eq:17.57}
\end{equation}
为把截面 \eqref{eq:17.49} 转换为末态部分子可观测量，我们需要 Jacobian
\begin{equation}
  \frac{\partial(x_1,x_2,\hat{t})}{\partial(y_3,y_4,p_\perp^2)} .
  \label{eq:17.58}
\end{equation}
把式~\eqref{eq:17.49} 乘以这个因子给出
\begin{equation}
  \frac{d\sigma}{dy_3\,dy_4\,dp_\perp^2}\bigl( p_1p_2 \to 3+4+X \bigr)
  = f_1(x_1)f_2(x_2)\ \frac{d\hat{\sigma}}{d\hat{t}}\bigl( 1+2 \to 3+4
    \bigr)\ \cdots .
  \label{eq:17.59}
\end{equation}
利用关系 $\hat{s} = x_1x_2s$ 和 $p_\perp dp_\perp = d^2p_\perp/2\pi$，这可以稍微化简，得到最终结果：
\begin{equation}
  \frac{d\sigma}{dy_3\,dy_4\,d^2p_\perp}\bigl( p_1p_2 \to 3+4+X \bigr)
  = f_1(x_1)f_2(x_2)\ \frac{d\hat{\sigma}}{d\hat{t}}\bigl( 1+2 \to 3+4
    \bigr)\ \frac{1}{\hat{s}}\ \cdots .
  \label{eq:17.60}
\end{equation}
在这个公式中，$x_1$、$x_2$ 以及部分子子过程的 Mandelstam 变量由式~\eqref{eq:17.57} 和 \eqref{eq:17.56} 给出。

这一结果给了我们任意二体部分子反应中末态轻子或喷注的完整分布。例如，为求 Drell-Yan 过程中末态轻子的分布，我们将把夸克湮灭为轻子的微分截面插入这个公式，
\begin{equation}
  \frac{d\hat{\sigma}}{d\hat{t}}\bigl( q\bar{q} \to \ell^{+}\ell^{-} \bigr)
  = \frac{2\pi\alpha^2 Q_f^2}{3\hat{s}^2}\ \frac{\hat{t}^2 + \hat{u}^2}
    {\hat{t}^2} .
  \label{eq:17.61}
\end{equation}
如果我们知道相关的部分子层次微分截面，这个公式同样可以应用于其他二体部分子反应。

\subsection{喷注对产生}

最常见的二体部分子反应是 QCD 的反应，涉及夸克、胶子或两者。遗憾的是，很难区分由胶子引发的强子喷注与由夸克引发的喷注。在实验上确定硬散射过程中的初始部分子是夸克还是胶子甚至更加困难。因此，QCD 对硬散射过程的预言最常以强子对撞中喷注产生的截面形式引用，对夸克、反夸克和胶子的所有可能反应求和。无论如何，为推导这些预言，我们必须算出基本的部分子-部分子截面。

夸克、反夸克和胶子的简单二体散射过程是 QCD 微扰论的基本过程，正如第~5 章研究的反应是 QED 微扰论的基本过程一样。它们是 QCD 中在 $\alpha_s$ 领头阶出现的基本强子硬散射反应。在本节的其余部分，我们将写下各种可能的夸克和胶子子过程的截面公式。所有这些截面都是 $\alpha_s^2$ 阶的。在实践中，这个 $\alpha_s$ 应该在反应的一个典型动量转移处求值，例如在 $Q^2 = \hat{t}$ 处。

\begin{figure}[t]
  \centering
  \includegraphics[width=0.45\textwidth]{fig17.9.png}
  \caption{对 $ud\to ud$ 有贡献的费曼图。}
  \label{fig:17.9}
\end{figure}

最简单的子过程是不同种类夸克的散射，例如 $u + d \to u + d$。在 $\alpha_s^2$ 阶，这一过程通过图~\ref{fig:17.9} 所示的费曼图发生。这一过程类似于 QED 中的电子-缪子散射，我们在式~\eqref{eq:17.22} 中写过它的截面：
\begin{equation}
  \frac{d\hat{\sigma}}{d\hat{t}}\bigl( e\mu \to e\mu \bigr)
  = \frac{2\pi\alpha^2}{\hat{s}^2}\ \frac{\hat{s}^2 + \hat{u}^2}{\hat{t}^2}
  .
  \label{eq:17.62}
\end{equation}
为把它转换为 QCD 中夸克散射的截面，我们只需把 QED 耦合 $e^2$ 换成 $g^2$ 乘以一个 $SU(3)$ 群论因子。QCD 图含有因子
\begin{equation}
  t^a_{ii'}\,t^a_{jj'} ,
\end{equation}
其中 $i$、$i'$ 是 $u$ 夸克的初末色，$j$、$j'$ 是 $d$ 夸克的初末色。为计算截面，我们必须对这个因子取平方，对末态色求和，并对初态色求平均。这给出因子
\begin{equation}
  \frac{1}{9}\ \mathrm{tr}\bigl[ t^at^b \bigr]\mathrm{tr}\bigl[ t^at^b \bigr]
  = \frac{1}{9}\ \bigl[ C(r) \bigr]^2\,\delta^{ab}\delta^{ab}
  = \frac{1}{9}\cdot\frac{1}{4}\cdot 8
  = \frac{2}{9} ,
  \label{eq:17.63}
\end{equation}
其中我们用式~\eqref{eq:15.78} 以及 $SU(3)$ 基础表示的 $C(r) = 1/2$。因此对 $ud$ 散射，
\begin{equation}
  \frac{d\hat{\sigma}}{d\hat{t}}\bigl( ud \to ud \bigr)
  = \frac{4\pi\alpha_s^2}{9\hat{s}^2}\ \frac{\hat{s}^2 + \hat{u}^2}
    {\hat{t}^2} .
  \label{eq:17.64}
\end{equation}
相同公式适用于任意两个不同夸克的散射，或通过交叉对称适用于一个夸克与一个不同种类的反夸克的散射。从 $t$ 道到 $s$ 道的交叉给出 $q\bar{q}$ 湮灭为不同种类的截面：
\begin{equation}
  \frac{d\hat{\sigma}}{d\hat{t}}\bigl( u\bar{u} \to d\bar{d} \bigr)
  = \frac{4\pi\alpha_s^2}{9\hat{s}^2}\ \frac{\hat{t}^2 + \hat{u}^2}
    {\hat{s}^2} .
  \label{eq:17.65}
\end{equation}

\begin{figure}[t]
  \centering
  \includegraphics[width=0.5\textwidth]{fig17.10.png}
  \caption{对 $u\bar{u}\to u\bar{u}$ 有贡献的费曼图。}
  \label{fig:17.10}
\end{figure}

夸克与同种反夸克的散射更加复杂，因为现在有两个相互干涉的费曼图，如图~\ref{fig:17.10} 所示。类似的 QED 过程是 Bhabha 散射 $e^{+}e^{-}\to e^{+}e^{-}$，我们在习题~5.2 中算出了它的截面：
\begin{equation}
  \frac{d\sigma}{dt}\bigl( e^{+}e^{-} \to e^{+}e^{-} \bigr)
  = \frac{2\pi\alpha^2}{s^2}\ \Bigl( \frac{s^2 + u^2}{t^2}
    + \frac{s^2 + t^2}{u^2} + \frac{4s^2}{tu} \Bigr) .
  \label{eq:17.66}
\end{equation}
然而，把它转写到 QCD 并非完全直截了当，因为不同的项得到不同的色因子。$t$ 道、$u$ 道和干涉项的色因子各不相同，必须分别计算。完成这一计算，人们得到
\begin{equation}
  \frac{d\hat{\sigma}}{d\hat{t}}\bigl( u\bar{u} \to u\bar{u} \bigr)
  = \frac{4\pi\alpha_s^2}{9\hat{s}^2}\ \Bigl( \frac{\hat{s}^2 + \hat{u}^2}
    {\hat{t}^2} + \frac{\hat{s}^2 + \hat{t}^2}{\hat{u}^2}
    - \frac{2}{3}\frac{\hat{s}^2}{\hat{t}\hat{u}} \Bigr) .
  \label{eq:17.67}
\end{equation}

其余的部分子子过程涉及胶子。夸克-胶子散射和胶子-胶子散射的截面在习题~17.3 中计算。结果为：
\begin{equation}
  \frac{d\hat{\sigma}}{d\hat{t}}\bigl( qg \to qg \bigr)
  = \frac{4\pi\alpha_s^2}{9\hat{s}^2}\ \Bigl( \frac{\hat{s}^2 + \hat{u}^2}
    {\hat{t}^2} - \frac{2}{3}\frac{\hat{s}}{\hat{u}} - \frac{2}{3}
      \frac{\hat{u}}{\hat{s}} + \frac{\hat{s}^2 + \hat{t}^2}{\hat{u}^2}
    \Bigr) ,
  \label{eq:17.68}
\end{equation}
\begin{equation}
  \frac{d\hat{\sigma}}{d\hat{t}}\bigl( gg \to gg \bigr)
  = \frac{4\pi\alpha_s^2}{\hat{s}^2}\ \frac{3}{2}\ \Bigl( \frac{\hat{s}^2
    + \hat{u}^2}{\hat{t}^2} + \frac{\hat{s}^2 + \hat{t}^2}{\hat{u}^2}
    - \frac{4}{3}\frac{\hat{s}^2}{\hat{t}\hat{u}} \Bigr) .
  \label{eq:17.69}
\end{equation}

利用这些部分子层次截面以及部分子模型公式 \eqref{eq:17.40}，我们可以计算强子对撞中双喷注产生的截面。图~\ref{fig:17.13} 把双喷注产生的领头阶 QCD 预言与高能 $p\bar{p}$ 对撞中观测到的双喷注事件不变质量分布作了比较。

\begin{figure}[t]
  \centering
  \includegraphics[width=0.6\textwidth]{fig17.13.png}
  \caption{$E_{cm} = 1.8$ TeV 的 $p\bar{p}$ 对撞中双喷注不变质量分布，由 CDF 合作组测量，F.~Abe 等人，\emph{Phys.\ Rev.\ D} \textbf{48}, 998 (1993)。该测量与使用图~\ref{fig:17.6} 所述 CTEQ 结构函数的领头阶 QCD 计算作了比较。三条较低曲线显示理论预言三个分量的不变质量分布：夸克-夸克（及反夸克）散射、夸克-胶子散射和胶子-胶子散射。}
  \label{fig:17.13}
\end{figure}

由于在部分子截面中求值 $\alpha_s(Q^2)$ 所用 $Q^2$ 的选择存在歧义，以及由于从深度非弹性散射截面推导部分子分布时存在类似歧义，理论预言的总体归一化不确定约一个因子 2。当计入 $\alpha_s$ 阶修正时，这一不确定性减小到约 30\%。尽管如此，值得注意的是，最低阶 QCD 预言跟随着观测到的双喷注不变质量分布，而后者下降六个数量级。因此，对于喷注产生截面，正如对涉及轻子的硬过程一样，QCD 确实对大动量转移下强相互作用的行为给出了合理的描述。

\section{部分子演化}

既然我们已经考察了 QCD 对几个强相互作用过程的领头阶预言，我们就应该研究这些预言在 $\alpha_s$ 下一阶的修正。我们在第~17.2 节看到，单个图的修正可能含有质量奇异性，即与共线发射过程相联系的、在零质量极限下出现的奇异性。对 $e^{+}e^{-}$ 湮灭为强子的过程，我们看到这些质量奇异性以及来自软胶子发射的红外发散在\emph{总}截面的表达式中抵消。可以证明，这是轻子或光子碰撞中产生夸克和胶子的过程的普遍特征。然而，当夸克或胶子出现在部分子子过程的\emph{初}态时，对过程的修正一般来说将具有不抵消的质量奇异性。在本节中，我们将演示这一效应并阐明它的物理解释。我们将发现，这些奇异项预言了 Bjorken 标度被依赖动量标度对数的项所破坏。事实上，它们导致一组支配部分子分布动量依赖的精确微分方程。

与 QCD 中质量奇异性相联系的基本现象已经存在于高能 QED 共线光子发射的物理中，因此从研究这种情形开始是最直截了当的。在本节中，我们将证明共线光子发射导致电子的一种部分子分布函数类比物。我们将推导描述这一分布函数的微分方程，它最早由 Gribov 和 Lipatov 构造。最后，我们将遵循 Altarelli 和 Parisi 的构造把这个方程推广到 QCD。%
\footnote{V.~N.~Gribov 和 L.~Lipatov，\emph{Sov.\ J.\ Nucl.\ Phys.} \textbf{15}, 438 (1972)；G.~Altarelli 和 G.~Parisi，\emph{Nucl.\ Phys.\ B} \textbf{126}, 298 (1977)。我们还强烈推荐阅读 J.~Kogut 和 L.~Susskind 的论文，\emph{Phys.\ Rev.\ D} \textbf{9}, 697、3391 (1974)。}

在第~5 章和第~6 章，我们研究了若干涉及含 $t$ 道或 $u$ 道奇异性图的 QED 过程例子。在这些情形下，我们发现总截面在高能极限下一般被一个额外因子 $\log(s/m^2)$ 增强。例如，在式~\eqref{eq:5.95} 中我们看到，Compton 散射中的 $u$ 道交换图（图~\ref{fig:17.14}(a)）导致一个在高能极限下取如下形式的积分
\begin{equation}
  \int d\cos\theta\ \frac{1}{(1+\cos\theta)} .
\end{equation}
$\cos\theta\to -1$ 处的奇异性被电子质量截断，导致对数增强因子。因此共线光子发射的代价因子不是 $\alpha$ 而是 $\alpha\log(s/m^2)$。如图~\ref{fig:17.14}(b) 所示的多重共线光子发射给出量级为 $(\alpha\log(s/m^2))^n$ 的贡献。为提高微扰论的精度，找到一种把这些项对 $\alpha$ 的所有阶求和的程序将是有用的。在 QCD 中，共线胶子发射的相应因子将是
\begin{equation}
  \alpha_s(Q^2)\log\frac{Q^2}{\mu^2} ,
\end{equation}
其中 $\mu$ 是非微扰 QCD 效应变得重要的动量标度。与式~\eqref{eq:17.17} 比较，我们看到这个乘积的量级为 1。因此，在这种情况下，如果我们要作出任何定量预言，大对数项的重求和是必不可少的。

\begin{figure}[t]
  \centering
  \includegraphics[width=0.55\textwidth]{fig17.14.png}
  \caption{与共线光子发射相联系并含质量奇异性的图：(a) 领头阶；(b) 高阶。}
  \label{fig:17.14}
\end{figure}

\begin{figure}[t]
  \centering
  \includegraphics[width=0.55\textwidth]{fig17.15.png}
  \caption{QED 中带质量奇异性图的一般形式。}
  \label{fig:17.15}
\end{figure}

在 QED 中，与一次共线发射相联系并含质量奇异性的图具有图~\ref{fig:17.15} 所示的其中一种形式。在每种情形中，圆圈代表大动量转移的散射过程。当中间传播子的分母为零时，即当中间态几乎在壳时，质量奇异性出现。因此，把图~\ref{fig:17.15} 中的第一个图看成到实光子和几乎实电子的跃迁，随后电子与振幅中其余粒子相互作用，是很自然的。第二个图应该有类似的解释，只是中间态是几乎实的光子。唯一的微妙之处在于定义中间态粒子的极化。对中间态为电子的情形，传播子的分子为
\begin{equation}
  \slashed{k} + m .
  \label{eq:17.79}
\end{equation}
因此，当 $k^2\to 0$ 时，光子发射顶点和振幅的其余部分与无质量电子的在壳极化旋量缩并。对中间态为光子的图的类似陈述是，电子发射顶点和剩余的光子振幅应该与中间态光子的物理横向极化矢量缩并。由于光子传播子的分子是 $g^{\mu\nu}$，光子传播子按这种方式约化并非显然。但这是正确的。为看清这一点，利用式~(16.20) 给出的 $g^{\mu\nu}$ 按无质量极化矢量的展开：
\begin{equation}
  g^{\mu\nu} = \epsilon^{\mu}_{+}\epsilon^{\nu*}_{-} + \epsilon^{\mu}_{-}
    \epsilon^{\nu*}_{+} - g^{\mu\nu}\ \cdots .
  \label{eq:17.80}
\end{equation}
这里 $\epsilon^{\mu}_{\pm}$ 是横向极化矢量。前向极化矢量 $\epsilon_{+}$ 正比于光子动量 $q^\mu$。当我们把 $\epsilon_{+}$ 与右侧的 QED 散射振幅缩并时，由 Ward 恒等式将得到零，$\epsilon_{+}$ 与电子发射顶点的缩并同样给出零。因此，为了在光子动量 $q$ 趋于在壳时计算奇异项，我们可以替换
\begin{equation}
  \frac{-ig^{\mu\nu}}{q^2} \to \frac{i}{q^2}\ \sum_{\lambda=\pm}
    \epsilon^{\mu*}_{\lambda}\epsilon^{\nu}_{\lambda} ,
  \label{eq:17.81}
\end{equation}
并用横向极化矢量求值光子发射和吸收振幅。

\subsection{电子分裂的矩阵元}

通过把中间传播子的分子替换为对极化矢量的求和，我们把光子或电子发射顶点与图的其余部分解耦。我们现在将在无质量粒子的物理极化态之间显式求值这个顶点。运动学如图~17.16 所示。两个末态粒子应该几乎共线，具有小的相对横向动量。我们可以选择入射电子动量沿 3 轴，出射动量在 1-3 平面内。设 $z$ 是初始电子能量中被光子带走的部分。于是三个 4-动量可以写成
\begin{equation}
  \begin{aligned}
    p &= (p, 0, 0, p),\\
    q &\approx (zp, p_\perp, 0, zp),\\
    k &\approx ((1-z)p, -p_\perp, 0, (1-z)p) .
  \end{aligned}
  \label{eq:17.82}
\end{equation}
这三个矢量满足 $p^2 = q^2 = k^2 = 0$，直至 $p_\perp^2$ 量级的项。在发射实光子的过程中，我们应该有 $p^2$ 和 $q^2$ 精确为零，而 $k^2$ 稍微离壳，离壳量级为 $p_\perp^2$。我们将需要知道出现在虚电子传播子中的 $k^2$ 的值。因此让我们修改式~\eqref{eq:17.82} 以满足 $q^2 = 0$ 直至 $p_\perp^2$ 量级项的条件，把 $q$ 和 $k$ 改写成
\begin{equation}
  \begin{aligned}
    q &= \Bigl( zp, p_\perp, 0, zp - \frac{p_\perp^2}{2p} \Bigr),\\
    k &= \Bigl( (1-z)p, -p_\perp, 0, (1-z)p + \frac{p_\perp^2}{2p}
      \Bigr) .
  \end{aligned}
  \label{eq:17.83}
\end{equation}
有了这一修改，
\begin{equation}
  k^2 = -\frac{p_\perp^2}{1-z}\ -\ 2(1-z)\,\cdots = -\frac{p_\perp^2}{1-z}
    + \mathcal{O}(p_\perp^4) .
  \label{eq:17.84}
\end{equation}
因此，如果光子是实的而电子是虚的，我们有
\begin{equation}
  q^2 = 0,\qquad k^2 = -\frac{p_\perp^2}{1-z} .
\end{equation}
反过来，在电子实而光子虚的过程中，
\begin{equation}
  k^2 = 0,\qquad q^2 = -\frac{p_\perp^2}{z} .
  \label{eq:17.85}
\end{equation}
这些更精确的表达式只在虚粒子的传播子中才需要。电子-光子顶点的矩阵元从 $p_\perp$ 阶开始，因此不受把 \eqref{eq:17.82} 修改为 \eqref{eq:17.83} 的显著影响，而且无论哪个粒子是虚的它都相同（到最低阶）。

我们现在计算 QED 顶点在具有确定螺旋度的无质量态之间的矩阵元。如果初始电子是左手的，由螺旋度守恒，末态电子也必须是左手的。于是光子发射顶点由
\begin{equation}
  i\mathcal{M} = \bar{u}_L(k)\,(-ie\gamma^\mu)\,u_L(p)\,
    \epsilon^{*}_\mu(q) ,
  \label{eq:17.86}
\end{equation}
给出，其中光子极化矢量既可以是左手的也可以是右手的。回顾螺旋度基表达式
\begin{equation}
  u_L(p) = \begin{pmatrix}\sqrt{p\cdot\sigma}\,\xi_L\\ \sqrt{p\cdot
    \bar{\sigma}}\,\xi_L\end{pmatrix}
  \qquad\text{对 } m = 0 ,
\end{equation}
我们可以更显式地写成
\begin{equation}
  i\mathcal{M} = ie\sqrt{2(1-z)p}\,\sqrt{2p}\ \xi^{\dagger}_L(k)\,
    \bm{\sigma}\cdot\bm{\epsilon}^{*}(q)\,\xi_L(p) .
  \label{eq:17.87}
\end{equation}
求值旋量积，我们发现只有当光子与电子具有相同螺旋度时振幅才非零（对左手电子发射光子的情形，由角动量守恒，光子必须是右手的）。对极化取平方并求和后，结果给出分裂概率
\begin{equation}
  P_{e\to e\gamma}(z) = \frac{\alpha}{2\pi}\ \frac{1+z^2}{1-z} .
\end{equation}

遵循导出式~(17.97) 的步骤，我们发现电子发射共线光子然后散射的过程的截面为
\begin{equation}
  \begin{aligned}
    \sigma(e^{-}\to Y) &= \int_0^1 dz\ \int \frac{dp_\perp^2}{p_\perp^2}\
      \frac{\alpha}{2\pi}\ \frac{1+(1-z)^2}{z}\ \hat{\sigma}(e^{-} \to Y)
      \ \cdots\\
    &= \int dz\ \int \frac{dp_\perp^2}{p_\perp^2}\ \frac{\alpha}{2\pi}\
      \frac{1+(1-z)^2}{z}\ \hat{\sigma}(e^{-} \to Y) ,
  \end{aligned}
  \label{eq:17.99}
\end{equation}
其中中间电子携带纵分数 $(1-z)$。

人们很自然会代入 $x = (1-z)$，并把 \eqref{eq:17.99} 中积分号下乘以截面的因子解释为在电子中找到电子部分子的部分子分布。这将给出
\begin{equation}
  f^{(1)}_e(x) = \frac{\alpha}{2\pi}\log\frac{Q^2}{m^2}\ \frac{1+x^2}{1-x}
  .
  \label{eq:17.100}
\end{equation}
然而，这个表达式并不充分。最明显的是，它没有考虑无辐射的过程，在这些过程中电子以全能量仍然是电子。这很容易补救：把 \eqref{eq:17.100} 看作对最朴素预期
\begin{equation}
  f^{(0)}_e(x) = \delta(1-x) ,
  \label{eq:17.101}
\end{equation}
的 $\alpha$ 阶修正即可，在该预期中我们认为电子只含有一个全能量的电子。遗憾的是，\eqref{eq:17.101} 和 \eqref{eq:17.100} 之和仍然不能给出电子分布的充分描述，原因有二。第一，式~\eqref{eq:17.100} 在 $x = 1$ 附近发散，我们需要一个处理这种奇异性的方案。第二，虽然式~\eqref{eq:17.100} 考虑了通过辐射从 $x = 1$ 移到纵分数 $x$ 的虚电子，它却没有考虑 $x = 1$ 处 delta 函数峰中电子随之而来的损失。

\eqref{eq:17.100} 在 $x = 1$ 处的发散对应于软光子的发射。我们在第~6.5 节看到，软光子的发射不影响 QED 反应的速率。对 $\alpha$ 逐阶地，人们发现来自软光子发射的对总速率的红外发散正贡献被来自含软虚光子图的负贡献所平衡。在当前例子中，负贡献必须减小不发射光子过程的权重。因此，到 $\alpha$ 阶，电子中电子的部分子分布应该具有如下形式
\begin{equation}
  f_e(x) = \delta(1-x) + \frac{\alpha}{2\pi}\log\frac{Q^2}{m^2}\ \Bigl[
    \frac{1+x^2}{(1-x)_+} - A\,\delta(1-x) \Bigr] .
  \label{eq:17.102}
\end{equation}
系数 $A$ 来自我们尚未计算的含虚光子的图。然而，这些图的效果容易理解；它们从 delta 函数中减去通过辐射移到较低 $x$ 的概率，因此对 $\alpha$ 阶完整项的积分为零。表达这一判据的另一种方式是，$A$ 由电子恰好含有一个电子部分子这个条件确定，
\begin{equation}
  \int_0^1 dx\ f_e(x) = 1 .
  \label{eq:17.103}
\end{equation}
（这个方程将在下面包含对产生过程时被修改。）

如何对 \eqref{eq:17.100} 中的奇异分母积分以显式确定 $A$ 并不那么清楚。通常的做法是定义一个可以通过从奇异项减去 delta 函数来积分的分布。定义分布
\begin{equation}
  \frac{1}{(1-x)_+}
  \label{eq:17.104}
\end{equation}
使其对所有小于 1 的 $x$ 值与函数 $1/(1-x)$ 一致，并在 $x = 1$ 处具有一个奇异性，使得该分布与任意光滑函数 $f(x)$ 的积分给出
\begin{equation}
  \int_0^1 dx\ \frac{f(x)}{(1-x)_+} = \int_0^1 dx\ \frac{f(x) - f(1)}
    {1-x} .
  \label{eq:17.105}
\end{equation}
不太形式化地，
\begin{equation}
  \frac{1}{(1-x)_+} = \lim_{\epsilon\to 0}\ \Bigl[ \frac{1}{1-x}
    \theta(1-\epsilon-x) - \delta(1-x)\ \int_0^{1-\epsilon} \frac{d\xi}{1-\xi}
    \Bigr] .
  \label{eq:17.106}
\end{equation}
更形式的定义 \eqref{eq:17.105} 在实践中往往更容易使用。

利用这个定义，我们可以通过把分母 $(1-x)$ 换成 $(1-x)_+$，把 delta 函数的一部分并入 \eqref{eq:17.102} 的奇异项。然后，为归一化 \eqref{eq:17.102}，我们需要积分
\begin{equation}
  \int_0^1 dx\ \frac{1+x^2}{1-x} = \int_0^1 dx\ \Bigl( \frac{x^2-1}{1-x}
    + \frac{2}{1-x} \Bigr) .
\end{equation}
到 $\alpha$ 阶，电子分布的最终形式为
\begin{equation}
  f_e(x) = \delta(1-x) + \frac{\alpha}{2\pi}\log\frac{Q^2}{m^2}\ \Bigl[
    \frac{1+x^2}{(1-x)_+} - \frac{3}{2}\delta(1-x) \Bigr] .
  \label{eq:17.107}
\end{equation}
这个分布现在被适当地归一化了，但它在 $x = 1$ 附近仍然高度奇异。因此，我们应该预期电子分布函数的高阶修正在这个区域很重要。于是，我们必须思考如何处理许多共线光子的发射。

\subsection{多重分裂}

事实上，把我们刚刚完成的分析推广到计入许多共线光子的发射并不困难。考虑图~\ref{fig:17.17}(a) 所示的过程，其中光子 1 以横向动量 $p_{1\perp}$ 辐射，光子 2 以横向动量 $p_{2\perp}$ 辐射。光子 2 的发射可以像我们上面做的那样计算。如果 $p_{2\perp} \gg p_{1\perp}$，与 $p_{2\perp}$ 相比，第一个虚电子非常接近质量壳，因此在计算光子 1 的发射时我们可以忽略它的虚度。双光子发射给出量级为
\begin{equation}
  \alpha^2\ \log^2\frac{Q^2}{m^2} .
\end{equation}
的贡献。在相反的极限 $p_{2\perp} \ll p_{1\perp}$ 下，没有 $p_{2\perp}^2$ 量级的分母，因此我们得不到双对数。只有在 $p_{2\perp} \ll p_{1\perp}$ 的情形下，$\alpha^2$ 阶的贡献才能与 $\alpha$ 阶的贡献相竞争。

\begin{figure}[t]
  \centering
  \includegraphics[width=0.55\textwidth]{fig17.17.png}
  \caption{含共线光子发射的高阶图：(a) 两个共线光子；(b) 许多共线光子。}
  \label{fig:17.17}
\end{figure}

这个论证推广到任意多个共线光子的发射，图~\ref{fig:17.17}(b)。光子相空间中对应于顺序
\begin{equation}
  p_{1\perp} \gg p_{2\perp} \gg p_{3\perp} \gg \cdots
  \label{eq:17.108}
\end{equation}
的积分区域给出含有因子
\begin{equation}
  \frac{1}{n!}\ \alpha^n\ \log^n\frac{Q^2}{m^2} .
  \label{eq:17.109}
\end{equation}
的贡献。如果光子横向动量以任何其他方式排序，来自该区域的贡献在同一 $\alpha$ 阶至少少一个大对数的幂次。如果条件 \eqref{eq:17.108} 成立，那么当从图的外部走向硬碰撞时，虚电子动量越来越离质量壳。在这种情况下，电子动量被称为\emph{强有序}的。

这一组结论有一个有趣的物理解释。由于当我们深入图中时中间电子越来越虚，把当物理电子在相继更小的距离标度上被分析时，把它们解释为物理电子的组分是很自然的。$k^2 \sim p_\perp^2$ 的中间电子可以被认为当物理电子的波函数以分辨率 $\Delta r \sim (p_\perp)^{-1}$ 被探测时变得可见的电子组分。在这种图景中，在一个分辨率下看到的电子可以在更精细的标度下被分辨成一个更虚的电子和若干光子。

无论从计算费曼图的视角还是从电子结构这一更宏大的视角，把电子分裂成虚电子加光子的过程想象成电子组分横向动量的连续演化过程都是有用的。为在数学上描述这一过程，我们引入电子和光子分布函数对 $p_\perp$ 的显式依赖。我们定义函数 $f_\gamma(x, Q)$ 和 $f_e(x, Q)$ 来给出在物理电子中找到纵分数为 $x$ 的光子或电子的概率，同时计入横向动量 $p_\perp < Q$ 的共线光子发射。

如果 $Q$ 稍微增加到 $Q + \Delta Q$，我们必须计入 $f_e(x, Q)$ 中的电子组分辐射一个 $Q < p_\perp < Q + \Delta Q$ 的光子的可能性。电子分裂出一个带走其能量分数 $z$ 的光子的微分概率为
\begin{equation}
  \frac{\alpha}{2\pi}\ \frac{dp_\perp^2}{p_\perp^2}\ \frac{1+(1-z)^2}{z} .
  \label{eq:17.110}
\end{equation}
因此新的光子分布可以如下计算：
\begin{equation}
  \begin{aligned}
    f_\gamma(x, Q + \Delta Q) &= f_\gamma(x, Q) + \frac{\alpha}{2\pi}\
      \frac{\Delta Q^2}{Q^2}\ \int_0^1 dx'\,dz\ \frac{1+(1-z)^2}{z}\,
      f_e(x', p_\perp)\,\delta(x - zx')\\
    &= f_\gamma(x, Q) + \frac{\alpha}{2\pi}\ \frac{\Delta Q^2}{Q^2}\
      \int_x^1 \frac{dz}{z}\ \frac{1+(1-z)^2}{z}\ f_e\Bigl(
        \frac{x}{z}, Q \Bigr) .
  \end{aligned}
  \label{eq:17.111}
\end{equation}
过渡到连续演化，我们发现函数 $f_\gamma(x, Q)$ 由积分-微分方程
\begin{equation}
  \frac{df_\gamma(x, Q)}{d\log Q} = \frac{\alpha}{2\pi}\ \int_x^1 \frac{dz}{z}\
    \frac{1+(1-z)^2}{z}\ f_e\Bigl( \frac{x}{z}, Q \Bigr) .
  \label{eq:17.112}
\end{equation}
确定。类似地，物理电子中组分电子的分布将随 $Q$ 演化，反映由于光子辐射在较低 $x$ 值处电子的出现，以及这些电子在较高 $x$ 处的消失。式~\eqref{eq:17.107} 中方括号内的项对这两种效应给出了正确的核算。

\subsection{演化方程}

收集所有分裂过程，我们可以为 QED 中的电子、正电子和光子分布函数写出一组完整的耦合演化方程：
\begin{equation}
  \frac{df_f(x,Q)}{d\log Q} = \frac{\alpha}{2\pi}\ \int_x^1 \frac{dz}{z}\
    \sum_{f'} \bigl[ P_{f\to f'}(z)\,f_{f'}\Bigl( \frac{x}{z}, Q \Bigr)
    \bigr] ,
  \label{eq:17.120}
\end{equation}
其中分裂函数 $P_{f\to f'}(z)$ 由
\begin{equation}
  \begin{aligned}
    P_{e\to e\gamma}(z) &= \frac{1+z^2}{1-z},\\
    P_{e\to \gamma e}(z) &= 1 + (1-z)^2,\\
    P_{\gamma\to e\bar{e}}(z) &= z^2 + (1-z)^2,\\
    P_{e\to e}(z) &= \frac{3}{2}\,\delta(1-z) .
  \end{aligned}
  \label{eq:17.121}
\end{equation}
给出。为得到与给定动量转移 $Q$ 相关的电子的分布函数，我们应该在 $Q = m$ 处用初始条件
\begin{equation}
  f_e(x,Q) = \delta(1-x),\qquad f_{\bar{e}}(x,Q) = 0,\qquad
  f_\gamma(x,Q) = 0 ,
  \label{eq:17.122}
\end{equation}
积分这些方程。用不同的初始条件，相同方程给出物理正电子或光子的分布函数。这些方程的解如式~(17.115) 那样用来计算涉及由电子、正电子或光子诱导的、含大动量转移过程的截面。

演化方程 \eqref{eq:17.120} 的构造方式是守恒电子数和纵动量。因此，强子部分子分布所满足的基本求和规则 \eqref{eq:17.36} 和 \eqref{eq:17.39} 也适用于 QED 分布函数。具体地说，电子的分布函数含有一个净电子组分，
\begin{equation}
  \int_0^1 dx\ \bigl[ f_e(x,Q) - f_{\bar{e}}(x,Q) \bigr] = 1 ,
  \label{eq:17.123}
\end{equation}
并核算物理电子的总动量，
\begin{equation}
  \int_0^1 dx\ x\ \bigl[ f_e(x,Q) + f_{\bar{e}}(x,Q) + f_\gamma(x,Q) \bigr]
  = 1 .
  \label{eq:17.124}
\end{equation}
利用式~\eqref{eq:17.120} 显式验证这些积分的值不依赖 $Q$ 是一个有教益的练习。

\subsection{Altarelli-Parisi 方程}

如果我们在 QED 中遇到与共线光子发射相联系的质量奇异性，我们也必然在 QCD 中遇到与共线胶子和夸克发射相联系的质量奇异性。如果我们用无质量夸克和胶子计算第~17.3 节和第~17.4 节讨论的领头阶部分子截面的 $\alpha_s$ 阶修正，我们将发现这些修正项在对共线位形积分时发散。因此部分子模型表达式，至少在它们最简单的形式下，在 $\alpha_s$ 的次领头阶就已经失效。

然而，假设 QCD 的奇异性不比 QED 的更严重，前一节的考虑告诉我们如何处理这些奇异项。在 QED 中，我们发现把与质量奇异性相联系的大修正包含在部分子分布中而不是硬散射截面中是自然的。从这种观点看，奇异项提供了部分子分布作为动量标度对数函数的演化方程的核。动量转移为 $Q$ 的硬散射在量级为 $Q^{-1}$ 的距离上探测电子。当电子波函数被分辨到非常小的标度时，它表现为一个只携带总纵动量一部分的组分电子，加上若干组分光子和电子-正电子对。这些组分中携带电子总动量很大一部分的任何一个都可以引发硬散射过程。

完全相同的逻辑适用于 QCD 截面的计算。来自共线胶子或夸克发射区域的贡献应该与部分子分布函数相联系，而不是与硬散射截面相联系。如果我们作这种联系，我们将发现部分子分布不再独立于表征硬散射过程的动量 $Q$；相反，它们现在随 $Q$ 对数地演化。例如，深度非弹性散射的基本方程 \eqref{eq:17.30} 将变成
\begin{equation}
  \frac{d^2\sigma}{dx\,dQ^2}\bigl( e^{-}p \to e^{-}X \bigr)
  = \frac{2\pi\alpha^2}{Q^4}\ \Bigl[ 1 + (1-y)^2 \Bigr]\ \sum_f Q_f^2\,
    f_f(x, Q) ,
  \label{eq:17.125}
\end{equation}
因此 Bjorken 标度将被破坏。由于这种破坏只在 $Q^2$ 的对数标度上发生，它将是一个微妙的效应，而近似的 Bjorken 标度仍将是 QCD 的一个预言。但 Bjorken 标度的破坏是不可避免的，因为 QCD 是一个在所有动量标度下都具有自由度的量子场论。当我们在越来越短的距离上探测质子波函数时，我们激发高动量自由度，并把波函数分辨成越来越多的夸克、反夸克和胶子。

由式~\eqref{eq:17.120} 支配的 QED 部分子分布的演化由参数 $\alpha/\pi$ 表征，因此当 $Q$ 改变一个因子 10 时，部分子分布改变 $\sim 1\%$。在 QCD 中，支配演化速率的相应因子应该是 $\alpha_s(Q)/\pi$。因此，当 $Q$ 很小时，演化迅速，微扰论中更高阶的贡献很重要。最终，演化的初始条件由质子波函数在大距离标度下的形式决定，而这种形式无法用费曼图计算。另一方面，当 $Q$ 大时——实践中远高于 1 GeV——演化变慢，并由微扰论的领头阶主导。在那种情况下，QCD 微扰论对部分子分布演化的形式作出精确预言，这些预言可以对照实验检验。

为推导 QCD 中部分子分布的演化方程，我们可以使用上面为 QED 使用的相同技术和逻辑。有一个微妙之处：在 $q^2\to 0$ 极限下把胶子传播子约化为横向极化态（式~\eqref{eq:17.81}）这一事实不能像在 QED 中那样简单地证明。然而，在非阿贝尔情形下结果也是正确的。%
\footnote{参见，例如，J.~Collins 和 D.~Soper，载于 A.~Mueller 主编，\emph{Quantum Chromodynamics}（World Scientific, Singapore, 1991）。} 一旦这个技术点得到解决，共线发射的运动学就与 QED 中完全相同。因此我们得到与 QED 中形式相同的演化方程，只需把 $\alpha$ 换成 $\alpha_s$、插入适当的色因子、并核算三胶子顶点的效应。

\begin{figure}[t]
  \centering
  \includegraphics[width=0.6\textwidth]{fig17.20.png}
  \caption{对 QCD 中部分子演化有贡献的三个顶点。}
  \label{fig:17.20}
\end{figure}

QCD 中的共线发射过程涉及图~\ref{fig:17.20} 所示的三个顶点。其中，前两个具有与图~17.16 和 17.19 所示相同的 Lorentz 结构。除了耦合常数的强度之外，唯一的不同在于色指标。我们将像在前面的分析中处理自旋那样处理色：对初态色求平均，对末态色求和。于是图~\ref{fig:17.20} 的第一个顶点（表示夸克分裂成夸克和胶子）得到色因子
\begin{equation}
  \frac{1}{2}\mathrm{tr}[t^at^a] = C_2(r) = \frac{4}{3} .
  \label{eq:17.126}
\end{equation}
第二个顶点（表示胶子分裂成夸克-反夸克对）得到因子
\begin{equation}
  \frac{1}{2}\mathrm{tr}[t^at^a] = \frac{1}{2} .
  \label{eq:17.127}
\end{equation}
图~\ref{fig:17.20} 中的第三个顶点表示胶子分裂成两个胶子，这是非阿贝尔情形中新的效应。通过取顶点在具有确定螺旋度的横向胶子态之间的矩阵元来计算这个顶点对演化方程的贡献是直截了当的。这一计算是习题~17.4 的主题。

通过核算所有这些效应，我们可以把 QED 演化方程 \eqref{eq:17.120} 修改成 QCD 中部分子分布的正确演化方程组。它们被称为\emph{Altarelli-Parisi 方程}。它们描述在标度 $Q$ 下可以被当作无质量的每种夸克和反夸克味的部分子分布 $f_q(x,Q)$、$f_{\bar{q}}(x,Q)$，连同胶子的部分子分布 $f_g(x,Q)$ 的耦合演化。显式地，
\begin{equation}
  \begin{aligned}
    \frac{df_g(x,Q)}{d\log Q} &= \frac{\alpha_s(Q^2)}{\pi}\ \int_x^1
      \frac{dz}{z}\ \Bigl[ P_{g\to g}(z)\,f_g\Bigl( \frac{x}{z}, Q \Bigr)
      + P_{g\to q}(z)\sum_f f_f\Bigl( \frac{x}{z}, Q \Bigr) \Bigr],\\
    \frac{df_f(x,Q)}{d\log Q} &= \frac{\alpha_s(Q^2)}{\pi}\ \int_x^1
      \frac{dz}{z}\ \Bigl[ P_{q\to q}(z)\,f_f\Bigl( \frac{x}{z}, Q \Bigr)
      + P_{q\to g}(z)\,f_g\Bigl( \frac{x}{z}, Q \Bigr) \Bigr] .
  \end{aligned}
  \label{eq:17.128}
\end{equation}
前三个分裂函数可以从式~\eqref{eq:17.121} 取来，乘以式~\eqref{eq:17.126} 和 \eqref{eq:17.127} 中计算的色因子：
\begin{equation}
  \begin{aligned}
    P_{q\to qg}(z) &= \frac{4}{3}\ \frac{1+z^2}{1-z},\\
    P_{q\to gq}(z) &= \frac{4}{3}\ \frac{1+(1-z)^2}{z},\\
    P_{g\to q\bar{q}}(z) &= \frac{1}{2}\ \bigl[ z^2 + (1-z)^2 \bigr] .
  \end{aligned}
  \label{eq:17.129}
\end{equation}
第四个分裂函数还需要习题~17.4 的计算；结果为
\begin{equation}
  P_{g\to g}(z) = 6\,\delta(1-z) + 6\ \frac{z}{(1-z)_+}
    + 6z(1-z) - \frac{2}{3}n_f\,\delta(1-z) .
  \label{eq:17.130}
\end{equation}
这个表达式中的最后一项正比于轻夸克味数 $n_f$，是与胶子分裂成 $q\bar{q}$ 对相联系的减除项。Altarelli-Parisi 方程描述任何强子或任何强子组分的部分子分布演化，直至不被大对数增强的 $\alpha_s$ 阶修正。

我们对 Altarelli-Parisi 方程的推导尊重 QCD 关于夸克数和纵动量的守恒定律。因此，方程必须尊重部分子模型求和规则 \eqref{eq:17.36} 和 \eqref{eq:17.39}。与 QED 情形一样，显式验证这些积分不依赖 $Q$ 是有教益的。

\begin{figure}[t]
  \centering
  \includegraphics[width=0.6\textwidth]{fig17.21.png}
  \caption{$Q = 2$、20 和 200 GeV 处 $u$ 夸克的部分子分布函数 $xf_u(x,Q)$，展示了根据 Altarelli-Parisi 方程的部分子演化效应。这些曲线取自图~\ref{fig:17.6} 所述 CTEQ 对深度非弹性散射数据的拟合。}
  \label{fig:17.21}
\end{figure}

在 QED 中，我们可以用演化方程显式计算电子的结构函数。在 QCD 中，这不再可能，因为积分方程所需的初始条件由 QCD 的强耦合区域决定，因此并非先验已知。然而，人们可以通过测量给定 $Q^2$ 值处深度非弹性散射的截面，在实验上确定质子结构的初始条件。然后人们可以预言较高 $Q^2$ 值处的结构函数，从而预言深度非弹性截面。这一分析有一个微妙之处：胶子分布不是在深度非弹性散射中直接测量的，但它确实进入夸克分布的演化方程。因此，关于深度非弹性散射 $Q^2$ 依赖的部分信息只是用于确定胶子分布。然而，胶子分布由动量求和规则 \eqref{eq:17.39} 绝对地归一化，因此即使这种分布必须从数据拟合，演化方程也具有预言力。

Altarelli-Parisi 方程预言部分子分布函数演化的一种特征形式，如图~\ref{fig:17.21} 所示。高 $x$ 处的部分子倾向于辐射并下降到较低的 $x$ 值。与此同时，新的部分子在低 $x$ 处作为这种辐射的产物而形成。因此，随着 $Q^2$ 增大，部分子分布在大 $x$ 处减小，在小 $x$ 处增长得快得多。我们可以把质子想象成具有越来越多的组分，但每个组分携带总动量的更小部分，随着它在越来越小的距离上被探测。

\section{$\alpha_s$ 的测量}

我们现在已经发展了微扰 QCD 的基本工具，并把它们应用于若干硬散射过程。在本节中，我们将回顾强耦合常数 $\alpha_s$ 的实验确定，它们检验 QCD 的定量预言。从各种测量中提取的 $\alpha_s$ 值汇总在表~\ref{tab:17.1} 中。表中显示了从这些测量中的每一个提取的 $\alpha_s$ 值，用参考约定中的值 $\alpha_s^{\overline{MS}}(m_Z^2)$ 表示。

\begin{table}[t]
  \centering
  \caption{从 QCD 实验得到的 $\alpha_s(m_Z)$ 值}
  \begin{tabular}{lll}
    \toprule
    过程 & $\alpha_s(m_Z)$ & $Q$ (GeV)\\
    \midrule
    深度非弹性散射 & 0.118 (6) & 1.7\\
    $\tau$ 轻子衰变中的 $R$ & 0.123 (4) & 1.8\\
    $J/\psi$、$\Upsilon$ 谱学 & 0.110 (6) & 2.3\\
    $W$ 产生的横向动量 & 0.121 (24) & 4.\\
    深度非弹性散射（演化） & 0.112 (4) & 5.\\
    $e^{+}e^{-}$ 湮灭中的事件形状 & 0.121 (6) & 5.8, 9.1\\
    $J/\psi$、$\Upsilon$ 衰变率 & 0.108 (10) & 9.5\\
    $e^{+}e^{-}$ 湮灭中的 $R$（20--65 GeV） & 0.124 (21) & 35.\\
    $Z^0$ 衰变中的 $R$ & 0.124 (7) & 91.2\\
    \bottomrule
  \end{tabular}
  \label{tab:17.1}
\end{table}

表中显示的 $\alpha_s(m_Z)$ 值是通过把实验结果拟合到用最小减除给出的微扰 QCD 理论表达式而得到的。$\alpha_s$ 的值已用重整化群方程演化到 $Q = m_Z$。$R$ 指到强子与到轻子的截面或分宽度之比。括号中的数字是最后显示数字位的标准误差。标为「$Q$」的列给出测量进行时 $Q$ 值的大致概念。（通常，这些测量对 $Q$ 的一个范围求平均，这种平均在引用的 $\alpha_s$ 值中已被考虑。）本表基于 I.~Hinchliffe 为粒子数据组撰写的文章中汇编的结果，\emph{Phys.\ Rev.\ D} \textbf{50}, 1297 (1994)。该文包含完整的参考文献集以及对这些确定中不确定性来源的讨论。

我们看到，有几个实验把 $\alpha_s$ 确定到 5\% 的精度，而且各种确定在这一水平上彼此一致。在图~\ref{fig:17.23} 中，我们绘制了表~\ref{tab:17.1} 中表示的、换算到共同标度之前的 $\alpha_s$ 原始值对获得每个值时所在动量标度 $Q$ 的图。这一比较给出 $\alpha_s$ 跑动的一个引人注目的直接验证。

\begin{figure}[t]
  \centering
  \includegraphics[width=0.6\textwidth]{fig17.23.png}
  \caption{对 $\alpha_s$ 的测量，按测量进行时的动量标度 $Q$ 绘制。此图通过把表~\ref{tab:17.1} 所列 $\alpha_s(m_Z)$ 的值演化回表中指示的 $Q$ 值而构造。$e^{+}e^{-}$ 事件形状的值已被拆分成对应 TRISTAN 和 LEP 加速器实验的两个点。这些值与初始条件为 $\alpha_s(m_Z) = 0.117$ 的重整化群演化理论预期作了比较。}
  \label{fig:17.23}
\end{figure}

在本章开头，我们只用几个简单原则就写下了一个强相互作用基本理论的候选者：夸克的存在及其量子数的辨认，以及夸克相互作用理论应该是渐近自由规范理论的想法。值得注意的是，这些简单的考虑引导我们得到了一个对强相互作用的描述，它在渐近自由可以用作计算工具的硬散射区中对广泛的现象定量正确。

\problems

\textbf{17.1} \emph{两圈重整化群关系。}
\begin{enumerate}[label=(\alph*)]
  \item 在微扰论的更高阶，QCD $\beta$ 函数的表达式将是一个级数
    \begin{equation}
      \beta(g) = -b_0\frac{g^3}{16\pi^2} - b_1\frac{g^5}{(16\pi^2)^2}
        + \cdots .
    \end{equation}
    对重整化群方程积分，并证明跑动耦合常数现在由
    \begin{equation}
      \alpha_s(Q) \approx \frac{4\pi}{b_0}\ \Bigl[ \frac{1}{\log(Q^2/\Lambda^2)}
        - \frac{b_1}{b_0}\ \frac{\log(\log(Q^2/\Lambda^2))}
          {(\log(Q^2/\Lambda^2))^2} + \cdots \Bigr],
    \end{equation}
    给出，其中被省略的项按 $(\log(Q^2/\Lambda^2))^{-2}$ 减小。
  \item 把这个公式与 $e^{+}e^{-}$ 湮灭截面的微扰级数结合：
    \begin{equation}
      \sigma(e^{+}e^{-}\to\text{强子}) = \sigma_0\cdot 3\cdot
        \sum_f Q_f^2\ \Bigl[ 1 + a_1\alpha_s + a_2\alpha_s^2
          + \mathcal{O}(\alpha_s^3) \Bigr] .
    \end{equation}
    系数 $a_2$ 依赖定义 $\alpha_s$ 的重整化条件的细节。证明大 $s$ 时 $\sigma(s)$ 渐近行为的领头两项只依赖 $b_0$ 和 $b_1$，且独立于 $a_2$ 和 $b_2$。因此 QCD $\beta$ 函数的前两个系数独立于重整化方案。
\end{enumerate}

\textbf{17.2} \emph{对胶子自旋的直接检验。} 在本习题中，我们比较 QCD 的预言与一个夸克相互作用由标量玻色子中介的模型的预言。设标量胶子与夸克的耦合由
\begin{equation}
  \delta\mathcal{L} = g_S\bar{q}q\,S ,
\end{equation}
给出，并定义 $\alpha_g = g^2/4\pi$。
\begin{enumerate}[label=(\alph*)]
  \item 利用第~I 部分末项工程 (b) 和 (c) 部分描述的技术，计算 $e^{+}e^{-}\to q\bar{q}S$ 到微扰论领头阶的截面。这个截面依赖 $q$、$\bar{q}$ 和 $S$ 的能量，我们像式~\eqref{eq:17.18} 那样把它们表示为电子束能量的分数 $x_1$、$x_2$、$x_3$。证明
    \begin{equation}
      \frac{d^2\sigma}{dx_1\,dx_2}(e^{+}e^{-}\to q\bar{q}S)
      = \frac{4\pi\alpha^2}{3s}\cdot 3Q_f^2\cdot
        \frac{\alpha_g}{\pi}\ \frac{x_2}{(1-x_1)(1-x_3)^2}\ \cdots .
    \end{equation}
  \item 在实践中，很难在实验上区分夸克与胶子，因为这两种粒子都表现为强子喷注。因此，设 $x_a$ 为 $x_1$、$x_2$、$x_3$ 中最大者，$x_b$ 为第二大者，$x_c$ 为最小者。对各种可能性求和，导出 $d^2\sigma/dx_adx_b$ 的表达式——既在 QCD 中（用式~\eqref{eq:17.18}），也在标量胶子模型中。证明这些模型可以通过它们在 $(x_a, x_b)$ 平面中的分布来区分。
\end{enumerate}

\textbf{17.3} \emph{夸克-胶子和胶子-胶子散射。}
\begin{enumerate}[label=(\alph*)]
  \item 计算 QCD 中夸克-反夸克湮灭到 $\alpha_s$ 领头阶的微分截面
    \begin{equation}
      \frac{d\hat{\sigma}}{d\hat{t}}\bigl( q\bar{q} \to gg \bigr)
    \end{equation}
    这最容易通过计算确定夸克和胶子螺旋度态之间的振幅来完成。忽略所有质量。使用显式的极化矢量和旋量，例如，
    \begin{equation}
      \epsilon^\mu = \frac{1}{\sqrt{2}}(0, 1, i, 0)
    \end{equation}
    用于沿 $+3$ 方向运动的右手胶子。你只需要考虑横向极化的胶子。由螺旋度守恒，只有初态 $q_L\bar{q}_R$ 和 $q_R\bar{q}_L$ 可以贡献；由宇称，这两个态给出相同的截面。因此只需计算三个过程
    \begin{equation}
      q_L\bar{q}_R \to g_Rg_R;\qquad
      q_L\bar{q}_R \to g_Rg_L;\qquad
      q_L\bar{q}_R \to g_Lg_L .
    \end{equation}
    的振幅。事实上，由 $CP$ 不变性，第一个和第三个过程具有相等的截面。计算振幅后，对它们取平方，并与色因子恰当地组合以构造各种螺旋度截面。最后，把这些组合起来构成对初态自旋和色求平均的总截面。
  \item 计算胶子-胶子散射的微分截面
    \begin{equation}
      \frac{d\hat{\sigma}}{d\hat{t}}\bigl( gg \to gg \bigr)
    \end{equation}
    有 16 种可能的螺旋度组合，但其中许多通过宇称和交叉对称彼此相关。全部 16 种可以由
    \begin{equation}
      g_Rg_R \to g_Rg_R;\qquad
      g_Rg_R \to g_Rg_L;\qquad
      g_Rg_R \to g_Lg_L .
    \end{equation}
    的三个振幅构造。证明这三个振幅中的最后两个为零。第一个可以用 Jacobi 恒等式极大地化简。尘埃落定之后，16 个极化胶子散射截面中只有三个非零。组合这些以计算自旋和色平均的微分截面。
\end{enumerate}

\textbf{17.4} \emph{胶子分裂函数。} 计算 Altarelli-Parisi 方程的胶子分裂函数 \eqref{eq:17.130}。为完成这一计算，首先计算图~\ref{fig:17.20} 所示三胶子顶点在确定螺旋度胶子态之间的矩阵元。组合这些矩阵元以推导 $x < 1$ 区域的分裂函数。然后固定分裂函数在 $x = 1$ 处的奇异性，使该函数具有正确的整体归一化。

\textbf{17.5} \emph{重夸克的光产生。} 考虑质量 $M$、电荷 $Q$ 的重夸克对光产生过程 $\gamma + p \to Q\bar{Q} + X$。如果 $M$ 足够大，任何对这个过程有贡献的图都必须涉及大动量转移；因此微扰 QCD 分析应该适用。这个想法在实践上已经适用于 $c$ 夸克对的产生。算出到 QCD 领头阶的截面。选择对这个反应给出领头贡献的部分子子过程，并写出截面的部分子模型表达式。你需要计算相关的子过程截面，但这可以直接取第~5 章中的一个 QED 计算。然后利用这一结果写出 $\gamma$-质子散射截面的表达式。

\textbf{17.6} \emph{部分子分布函数在小 $x$ 处的行为。} 利用一些有物理动机的近似，有可能对很小的 $x$ 解析地求解 Altarelli-Parisi 方程。在这个极限下，胶子分布占主导，演化由 $P_{g\to g}(z)$ 在小 $z$ 处的奇异性驱动。证明当 $x\to 0$ 时胶子分布按 $x$ 的幂次增长，并用 $\alpha_s$ 推导出该幂次。
